% 子文件：4.多元函数积分学.tex
% 对应教材：同济《高等数学》下册 第十章 重积分
% 节次：§10.1 二重积分的概念与性质｜§10.2 二重积分的计算法（直角坐标·极坐标·换元法）
%       §10.3 三重积分（概念·直角坐标·柱面坐标·球面坐标）｜§10.4 重积分的应用｜*§10.5 含参变量的积分
% 主文档使用 % 子文件：4.多元函数积分学.tex
% 对应教材：同济《高等数学》下册 第十章 重积分
% 节次：§10.1 二重积分的概念与性质｜§10.2 二重积分的计算法（直角坐标·极坐标·换元法）
%       §10.3 三重积分（概念·直角坐标·柱面坐标·球面坐标）｜§10.4 重积分的应用｜*§10.5 含参变量的积分
% 主文档使用 % 子文件：4.多元函数积分学.tex
% 对应教材：同济《高等数学》下册 第十章 重积分
% 节次：§10.1 二重积分的概念与性质｜§10.2 二重积分的计算法（直角坐标·极坐标·换元法）
%       §10.3 三重积分（概念·直角坐标·柱面坐标·球面坐标）｜§10.4 重积分的应用｜*§10.5 含参变量的积分
% 主文档使用 % 子文件：4.多元函数积分学.tex
% 对应教材：同济《高等数学》下册 第十章 重积分
% 节次：§10.1 二重积分的概念与性质｜§10.2 二重积分的计算法（直角坐标·极坐标·换元法）
%       §10.3 三重积分（概念·直角坐标·柱面坐标·球面坐标）｜§10.4 重积分的应用｜*§10.5 含参变量的积分
% 主文档使用 \input{4.多元函数积分学} 引入，不含 document 环境

\section{二重积分的概念与性质}
\subsection{引例与定义}

\textbf{引例一：曲顶柱体的体积。}底为 $xOy$ 面上的闭区域 $D$，顶为曲面 $z=f(x,y)$（$f\ge0$ 连续），高随点而变，"体积 $=$ 高 $\times$ 底面积"失效。用曲线网把 $D$ 分成 $n$ 个小闭区域 $\Delta\sigma_i$（既表示小块，也表示其面积），细曲顶柱体近似为以 $f(\xi_i,\eta_i)$ 为高的平顶柱体，求和后令最大直径 $\lambda\to0$ 取极限：$V=\lim\limits_{\lambda\to0}\sum\limits_{i=1}^{n}f(\xi_i,\eta_i)\Delta\sigma_i$。

\textbf{引例二：平面薄片的质量。}薄片占有 $D$，面密度 $\mu(x,y)>0$ 连续，同样"分割—近似—求和—取极限"，得 $m=\lim\limits_{\lambda\to0}\sum\limits_{i=1}^{n}\mu(\xi_i,\eta_i)\Delta\sigma_i$。

两问题的实际意义不同，但所求量都归结为同一形式的和的极限，抽去具体意义即得

\textbf{定义（二重积分）}\quad 设 $f(x,y)$ 是有界闭区域 $D$ 上的有界函数。将 $D$ 任意分成 $n$ 个小闭区域 $\Delta\sigma_i$，在每个 $\Delta\sigma_i$ 上任取一点 $(\xi_i,\eta_i)$，作和 $\sum\limits_{i=1}^{n}f(\xi_i,\eta_i)\Delta\sigma_i$。若当 $\lambda\to0$ 时该和的极限总存在，且与 $D$ 的分法及点 $(\xi_i,\eta_i)$ 的取法无关，则称此极限为 $f(x,y)$ 在 $D$ 上的\textbf{二重积分}，记作
\[
\iint_D f(x,y)\,\mathrm{d}\sigma=\lim_{\lambda\to0}\sum_{i=1}^{n}f(\xi_i,\eta_i)\,\Delta\sigma_i ,
\]
其中 $f(x,y)$ 为\textbf{被积函数}，$f(x,y)\mathrm{d}\sigma$ 为\textbf{被积表达式}，$\mathrm{d}\sigma$ 为\textbf{面积元素}，$x,y$ 为\textbf{积分变量}，$D$ 为\textbf{积分区域}。

\begin{itemize}
  \item 直角坐标中用平行于坐标轴的直线网划分 $D$，除含边界点的小块外都是矩形，$\Delta\sigma=\Delta x\,\Delta y$，故 $\mathrm{d}\sigma=\mathrm{d}x\,\mathrm{d}y$，二重积分也记作 $\iint_Df(x,y)\,\mathrm{d}x\,\mathrm{d}y$。
  \item \textbf{可积性}：$f(x,y)$ 在 $D$ 上\textbf{连续}时积分和极限必存在，即必可积（可积的必要条件是被积函数有界）。教材以后总假定可积。
  \item \textbf{几何意义}：$f\ge0$ 时积分值是以 $D$ 为底、$z=f(x,y)$ 为顶的曲顶柱体体积；$f$ 变号时，积分值等于 $xOy$ 面上方柱体体积减去下方柱体体积之差。
\end{itemize}

\subsection{二重积分的性质}

设所涉积分都存在，$\sigma$ 为 $D$ 的面积。

\begin{enumerate}
  \item \textbf{线性性质}\quad $\iint_D[\alpha f+\beta g]\,\mathrm{d}\sigma=\alpha\iint_Df\,\mathrm{d}\sigma+\beta\iint_Dg\,\mathrm{d}\sigma$（$\alpha,\beta$ 为常数）。
  \item \textbf{对积分区域的可加性}\quad 若 $D$ 被有限条曲线分为有限个部分闭区域（如 $D=D_1\cup D_2$，无公共内点），则 $\iint_Df\,\mathrm{d}\sigma=\iint_{D_1}f\,\mathrm{d}\sigma+\iint_{D_2}f\,\mathrm{d}\sigma$。
  \item \textbf{面积性质}\quad 若在 $D$ 上 $f\equiv1$，则 $\iint_D1\,\mathrm{d}\sigma=\iint_D\mathrm{d}\sigma=\sigma$（高为 $1$ 的平顶柱体体积在数值上等于底面积）。
  \item \textbf{比较性质}\quad 若在 $D$ 上 $f\le g$，则 $\iint_Df\,\mathrm{d}\sigma\le\iint_Dg\,\mathrm{d}\sigma$；特别地 $\left|\iint_Df\,\mathrm{d}\sigma\right|\le\iint_D|f|\,\mathrm{d}\sigma$。
  \item \textbf{估值性质}\quad 设 $M,m$ 为 $f$ 在 $D$ 上的最大值、最小值，则 $m\sigma\le\iint_Df\,\mathrm{d}\sigma\le M\sigma$。
  \item \textbf{中值定理}\quad 设 $f$ 在 $D$ 上连续，则 $D$ 上至少存在一点 $(\xi,\eta)$ 使 $\iint_Df(x,y)\,\mathrm{d}\sigma=f(\xi,\eta)\sigma$。（证：由性质 5 得 $m\le\frac{1}{\sigma}\iint_Df\,\mathrm{d}\sigma\le M$，再用连续函数的介值定理。）
\end{enumerate}

\corepoint{性质 5、6 中的"$\sigma$ 是 $D$ 的面积"不能漏：直角坐标下 $\sigma=\iint_D\mathrm{d}x\mathrm{d}y$，极坐标下 $\sigma=\iint_D\rho\,\mathrm{d}\rho\mathrm{d}\theta$，换元后面积要重算，不能把 $\Delta x\Delta y$ 当面积。另外"$f\le g\Rightarrow$ 积分值 $\le$"只对\textbf{同一积分区域}成立。}

\subsection{对称性简化（常用）}

\begin{itemize}
  \item $D$ 关于 $y$ 轴对称时：
  \[
  \iint_Df(x,y)\,\mathrm{d}\sigma=
  \begin{cases}
  0, & f(-x,y)=-f(x,y),\\[2pt]
  2\displaystyle\iint_{D_1}f(x,y)\,\mathrm{d}\sigma, & f(-x,y)=f(x,y),
  \end{cases}
  \]
  $D_1$ 为 $D$ 中 $x\ge0$ 的部分。
  \item $D$ 关于 $x$ 轴对称时把 $x,y$ 互换；若 $D$ 关于两轴均对称且 $f$ 关于 $x,y$ 都是偶函数，则 $\iint_Df\,\mathrm{d}\sigma=4\iint_{D_1}f\,\mathrm{d}\sigma$（$D_1$ 为第一象限部分）；$D$ 关于原点对称且 $f(-x,-y)=f(x,y)$ 时 $\iint_Df\,\mathrm{d}\sigma=2\iint_{D_{y\ge0}}f\,\mathrm{d}\sigma$。
\end{itemize}

\corepoint{用对称性有\textbf{两个前提}：(1) 积分\textbf{区域}对称；(2) 被积函数在该对称下有确定的奇偶性，缺一不可。判断奇偶性时只有一个变量变号，其余当常数。若 $D$ 关于 $y$ 轴对称而 $f$ 关于 $x$ 为奇函数，则积分为 $0$，这一句常能省掉整段计算。}

\section{二重积分的计算法}

按定义计算只对少数特别简单的被积函数与区域可行。基本方法是把二重积分化为\textbf{两次定积分}（累次积分）。

\subsection{利用直角坐标计算}

取 $\mathrm{d}\sigma=\mathrm{d}x\,\mathrm{d}y$。

\textbf{$X$ 型区域（先对 $y$、后对 $x$）}\quad 若 $D$ 可用不等式 $a\le x\le b$，$\varphi_1(x)\le y\le\varphi_2(x)$ 表示（$\varphi_1\le\varphi_2$ 在 $[a,b]$ 上连续），则
\keyformula{\iint_Df(x,y)\,\mathrm{d}x\,\mathrm{d}y=\int_a^b\mathrm{d}x\int_{\varphi_1(x)}^{\varphi_2(x)}f(x,y)\,\mathrm{d}y}
特征：穿过 $D$ 内部且\textbf{平行于 $y$ 轴}的直线与 $D$ 的边界相交不多于两点。

\textbf{$Y$ 型区域（先对 $x$、后对 $y$）}\quad 若 $D$ 可用 $c\le y\le d$，$\psi_1(y)\le x\le\psi_2(y)$ 表示，则
\keyformula{\iint_Df(x,y)\,\mathrm{d}x\,\mathrm{d}y=\int_c^d\mathrm{d}y\int_{\psi_1(y)}^{\psi_2(y)}f(x,y)\,\mathrm{d}x}
特征：穿过 $D$ 内部且\textbf{平行于 $x$ 轴}的直线与边界相交不多于两点。

两种情形要处理：

\begin{itemize}
  \item 若既有平行于 $y$ 轴的直线与边界相交多于两点、又有平行于 $x$ 轴的直线相交多于两点，则 $D$ 既非 $X$ 型也非 $Y$ 型，须把 $D$ 分成几块，每块为 $X$ 型或 $Y$ 型，再由性质 2 相加。
  \item 若 $D$ 既是 $X$ 型又是 $Y$ 型，则两种次序的二次积分相等，都等于同一个二重积分（这正是累次积分可交换次序的依据）。
\end{itemize}

\textbf{化累次积分的步骤}

\begin{enumerate}
  \item \textbf{画区域}：画出 $D$，求出边界曲线的全部交点。
  \item \textbf{定类型}：判断 $X$ 型、$Y$ 型还是需分块。
  \item \textbf{定限}：外层积分限必为\textbf{常数}，内层积分限是外层变量的函数。
  \item \textbf{算内层}：把外层变量当常数，得到外层变量的函数。
  \item \textbf{算外层}：对剩下的变量作定积分。
\end{enumerate}

\corepoint{定限是唯一的关键，口诀"先画图、后定限，外层常数、内层函数"。常见错误：外层限写成函数；上下限写反；分块时漏块；先对 $y$ 积时把含 $x$ 的因子留在内层（应作为常数提出）。}

\corepoint{\textbf{选择次序}要同时看区域形状与被积函数：既要 $D$ 少分块，又要内层能积出来。含 $\mathrm{e}^{y^2}$、$\frac{\sin y}{y}$、$\sqrt{1+y^3}$ 之类的积分对 $y$ 积不出来，必须\textbf{交换次序}：先由原累次积分还原 $D$ 的不等式组，画图，再按新次序重新定限。这常常是唯一出路。}

\subsection{利用极坐标计算}

当 $D$ 的边界用极坐标表示方便（圆、圆环、扇形）且被积函数含 $x^2+y^2$、$\frac{y}{x}$ 时用极坐标。取 $x=\rho\cos\theta$，$y=\rho\sin\theta$，两族曲线围成的小块面积 $\approx\rho\,\mathrm{d}\rho\,\mathrm{d}\theta$，故\textbf{极坐标下的面积元素} $\mathrm{d}\sigma=\rho\,\mathrm{d}\rho\,\mathrm{d}\theta$，
\keyformula{\iint_Df(x,y)\,\mathrm{d}x\,\mathrm{d}y=\iint_Df(\rho\cos\theta,\rho\sin\theta)\,\rho\,\mathrm{d}\rho\,\mathrm{d}\theta}
即：$x,y$ 分别换成 $\rho\cos\theta,\rho\sin\theta$，$\mathrm{d}x\mathrm{d}y$ 换成 $\rho\,\mathrm{d}\rho\mathrm{d}\theta$（\textbf{这个 $\rho$ 不能漏}）。

设 $D$ 为 $\rho_1(\theta)\le\rho\le\rho_2(\theta)$，$\alpha\le\theta\le\beta$，则
\keyformula{\iint_Df(\rho\cos\theta,\rho\sin\theta)\,\rho\,\mathrm{d}\rho\,\mathrm{d}\theta=\int_\alpha^\beta\mathrm{d}\theta\int_{\rho_1(\theta)}^{\rho_2(\theta)}f(\rho\cos\theta,\rho\sin\theta)\,\rho\,\mathrm{d}\rho}
两种特例：极点在 $D$ 的\textbf{边界}上（曲边扇形）时 $\rho_1\equiv0$，$\rho$ 的上限为 $\rho(\theta)$；极点在 $D$ 的\textbf{内部}时 $0\le\rho\le\rho(\theta)$，$0\le\theta\le2\pi$。由性质 3 得曲边扇形的面积 $\sigma=\frac{1}{2}\int_\alpha^\beta[\rho_2^2(\theta)-\rho_1^2(\theta)]\,\mathrm{d}\theta$。

\corepoint{极坐标三个高频错误：(1) 漏掉面积元素里的 $\rho$；(2) 不管区域张多大角，一律写 $[0,2\pi]$；(3) 内外层限的位置搞反。另需记准两个常用转换：$x^2+y^2=2ax$ 即 $\rho=2a\cos\theta$（$\theta\in[-\frac{\pi}{2},\frac{\pi}{2}]$），$x^2+y^2=2ay$ 即 $\rho=2a\sin\theta$（$\theta\in[0,\pi]$）。}

\subsection{二重积分的换元法}

设 $f(x,y)$ 在 $D$ 上连续，变换 $x=x(u,v)$，$y=y(u,v)$ 把 $uOv$ 面上的闭区域 $D'$ 一一对应地映成 $D$，$x,y$ 在 $D'$ 上有一阶连续偏导数，其\textbf{雅可比行列式}
\[
J(u,v)=\frac{\partial(x,y)}{\partial(u,v)}
=\begin{vmatrix}
\dfrac{\partial x}{\partial u} & \dfrac{\partial x}{\partial v}\\[8pt]
\dfrac{\partial y}{\partial u} & \dfrac{\partial y}{\partial v}
\end{vmatrix}\ne0 ,
\]
则
\keyformula{\iint_Df(x,y)\,\mathrm{d}x\,\mathrm{d}y=\iint_{D'}f\bigl(x(u,v),y(u,v)\bigr)\,\bigl|J(u,v)\bigr|\,\mathrm{d}u\,\mathrm{d}v}
极坐标 $x=\rho\cos\theta,y=\rho\sin\theta$ 是其特例，$|J|=\rho$。常用\textbf{广义极坐标}（椭圆域 $\frac{x^2}{a^2}+\frac{y^2}{b^2}\le1$）：$x=a\rho\cos\theta$，$y=b\rho\sin\theta$，$|J|=ab\rho$。

\corepoint{公式里是 $|J|$ 的\textbf{绝对值}；换元后区域一并换成 $D'$，积分限须按 $D'$ 重定；变换须一一对应（$J\ne0$ 只保证局部一一），否则分块。}

\section{三重积分}
\subsection{概念}

\textbf{定义}\quad 设 $f(x,y,z)$ 是空间有界闭区域 $\Omega$ 上的有界函数。把 $\Omega$ 任意分成 $n$ 个小闭区域 $\Delta V_i$（也表示其体积），任取 $(\xi_i,\eta_i,\zeta_i)$ 作和 $\sum f(\xi_i,\eta_i,\zeta_i)\Delta V_i$。若 $\lambda\to0$（各小块直径最大值）时极限总存在且与分法、取法无关，则称之为 $f$ 在 $\Omega$ 上的\textbf{三重积分}：
\[
\iiint_\Omega f(x,y,z)\,\mathrm{d}V=\lim_{\lambda\to0}\sum_{i=1}^{n}f(\xi_i,\eta_i,\zeta_i)\,\Delta V_i ,
\]
$\mathrm{d}V$ 为\textbf{体积元素}；直角坐标下 $\mathrm{d}V=\mathrm{d}x\,\mathrm{d}y\,\mathrm{d}z$。$f$ 连续时必可积；$f\ge0$ 时可解释为空间物体的质量（$f$ 为体密度）。

性质与二重积分完全平行：线性、对区域可加、$\iiint_\Omega\mathrm{d}V=V$、比较、估值、中值定理 $\iiint_\Omega f\,\mathrm{d}V=f(\xi,\eta,\zeta)V$。

\subsection{直角坐标：投影法与截面法}

\textbf{投影法（先一后二）}\quad 设平行于 $z$ 轴且穿过 $\Omega$ 内部的直线与边界曲面 $S$ 相交不多于两点。把 $\Omega$ 投影到 $xOy$ 面得 $D_{xy}$，$S$ 被分成下、上两部分 $S_1:z=z_1(x,y)$，$S_2:z=z_2(x,y)$（$z_1\le z_2$ 均在 $D_{xy}$ 上连续），则
\keyformula{\iiint_\Omega f\,\mathrm{d}V=\iint_{D_{xy}}\mathrm{d}x\,\mathrm{d}y\int_{z_1(x,y)}^{z_2(x,y)}f(x,y,z)\,\mathrm{d}z}
即先对 $z$、再对 $x,y$ 的三次积分。若平行于 $x$ 轴或 $y$ 轴的直线与 $S$ 相交不多于两点，也可投影到 $yOz$ 面或 $zOx$ 面，化为其他顺序的三次积分；交点多于两点时把 $\Omega$ 分块相加。

\textbf{截面法（先二后一）}\quad 设 $\Omega=\{(x,y,z)\mid c\le z\le d,\ (x,y)\in D_z\}$，其中 $D_z$ 是用竖坐标 $z$ 的平面截 $\Omega$ 所得的平面闭区域，则
\keyformula{\iiint_\Omega f\,\mathrm{d}V=\int_c^d\mathrm{d}z\iint_{D_z}f(x,y,z)\,\mathrm{d}x\,\mathrm{d}y}
当被积函数只含 $z$（或含 $z$ 的因子可分离）且截面面积易求时，内层二重积分 $=$ 截面面积 $\times$ 该因子。

\corepoint{两法的差别只在"先算哪个"：投影法是\textbf{先一后二}，要求沿 $z$ 方向穿入穿出各一次；截面法是\textbf{先二后一}，要求截面形状随 $z$ 变化简单。判据：被积函数中 $z$ 的因子能否提出内层——能提就优先截面法。}

\subsection{柱面坐标}

设点 $M$ 在 $xOy$ 面上投影 $P$ 的极坐标为 $(\rho,\theta)$，则 $(\rho,\theta,z)$ 为 $M$ 的\textbf{柱面坐标}，$0\le\rho<+\infty$，$0\le\theta\le2\pi$，$-\infty<z<+\infty$。三组坐标面：$\rho=$ 常数（以 $z$ 轴为轴的圆柱面）、$\theta=$ 常数（过 $z$ 轴的半平面）、$z=$ 常数（平行于 $xOy$ 面的平面）。关系式与体积元素：
\keyformula{x=\rho\cos\theta,\quad y=\rho\sin\theta,\quad z=z;\qquad \mathrm{d}V=\rho\,\mathrm{d}\rho\,\mathrm{d}\theta\,\mathrm{d}z}
\keyformula{\iiint_\Omega f\,\mathrm{d}V=\iiint_{\Omega'}f(\rho\cos\theta,\rho\sin\theta,z)\,\rho\,\mathrm{d}\rho\,\mathrm{d}\theta\,\mathrm{d}z}
再化为三次积分，积分限按 $\rho,\theta,z$ 在 $\Omega$ 中的变化范围确定。

\corepoint{柱面坐标适合"绕 $z$ 轴旋转"的区域（圆柱、圆锥、旋转抛物面、圆环）与被积函数含 $x^2+y^2$ 的情形。$\rho$ 的限与二重积分的极坐标一致：先看投影区域张多大角，再看沿射线 $\rho$ 从哪进哪出；别忘了因子 $\rho$。}

\subsection{球面坐标}

设 $r$ 为原点 $O$ 与 $M$ 的距离，$\varphi$ 为有向线段 $\overrightarrow{OM}$ 与 $z$ 轴正向的夹角，$\theta$ 为从 $z$ 轴正向来看自 $x$ 轴按逆时针转到 $\overrightarrow{OP}$ 的角（$P$ 为 $M$ 在 $xOy$ 面上的投影），则 $(r,\varphi,\theta)$ 为 $M$ 的\textbf{球面坐标}，$0\le r<+\infty$，$0\le\varphi\le\pi$，$0\le\theta\le2\pi$。三组坐标面：$r=$ 常数（以原点为球心的球面）、$\varphi=$ 常数（以原点为顶点、$z$ 轴为轴的圆锥面）、$\theta=$ 常数（过 $z$ 轴的半平面）。
\keyformula{x=r\sin\varphi\cos\theta,\quad y=r\sin\varphi\sin\theta,\quad z=r\cos\varphi}
\keyformula{\mathrm{d}V=r^{2}\sin\varphi\,\mathrm{d}r\,\mathrm{d}\varphi\,\mathrm{d}\theta}
\keyformula{\iiint_\Omega f\,\mathrm{d}V=\iiint_{\Omega'}f(r\sin\varphi\cos\theta,\ r\sin\varphi\sin\theta,\ r\cos\varphi)\,r^{2}\sin\varphi\,\mathrm{d}r\,\mathrm{d}\varphi\,\mathrm{d}\theta}
一般化为先对 $r$、再对 $\varphi$、最后对 $\theta$ 的三次积分：$0\le\theta\le2\pi$，$0\le\varphi\le\pi$（只在上半球时 $0\le\varphi\le\frac{\pi}{2}$），$r$ 的限由"自原点出发的射线穿过 $\Omega$"确定。

\corepoint{两个最易错处：(1) 体积元素是 $r^2\sin\varphi$，不是 $\cos\varphi$、也不是 $\sin\theta$；(2) $\varphi$ 是\textbf{与 $z$ 轴正向的夹角}，不是与 $xOy$ 面的夹角——上半球对应 $\varphi\in[0,\frac{\pi}{2}]$，圆锥面 $z=\sqrt{x^2+y^2}$ 对应 $\varphi=\frac{\pi}{4}$。记不住就用 $z=r\cos\varphi$ 现推。}

\subsection{三种坐标的选择}

\begin{center}
\begin{tabular}{@{}lll@{}}
\toprule
坐标 & 适用的积分区域 & 体积元素 \\
\midrule
直角坐标 & 长方体、四面体、由平面与柱面围成的可分层区域 & $\mathrm{d}x\,\mathrm{d}y\,\mathrm{d}z$ \\
柱面坐标 & 圆柱、圆锥、旋转抛物面、绕 $z$ 轴的旋转体 & $\rho\,\mathrm{d}\rho\,\mathrm{d}\theta\,\mathrm{d}z$ \\
球面坐标 & 球体、球壳、球与锥面所围的"冰淇淋"形 & $r^{2}\sin\varphi\,\mathrm{d}r\,\mathrm{d}\varphi\,\mathrm{d}\theta$ \\
\bottomrule
\end{tabular}
\end{center}

选择原则：先看\textbf{边界曲面}（球用球面坐标、柱面用柱面坐标），再看\textbf{被积函数}含 $x^2+y^2$（柱面坐标）还是 $x^2+y^2+z^2$（球面坐标）；冲突时以"积分限好定"为准。

\section{重积分的应用}
\subsection{元素法}

若所求量 $U$ (1) 对积分区域具有可加性，(2) 在区域内某点的很小邻域上其部分量可用 $f(M)\,\mathrm{d}\Omega$ 近似（$\mathrm{d}\Omega$ 为面积或体积元素），则 $U=\iint_Df\,\mathrm{d}\sigma$ 或 $U=\iiint_\Omega f\,\mathrm{d}V$。关键是写出\textbf{正确的微元}：几何量用面积（体积）元素，物理量用"密度 $\times$ 元素"。

\subsection{曲面的面积}

设曲面 $S:z=f(x,y)$，$f$ 在 $D$（$S$ 在 $xOy$ 面上的投影区域）有连续偏导。$S$ 上小片面积 $\mathrm{d}A$ 与其投影 $\mathrm{d}\sigma$ 之比为 $\frac{1}{\cos\gamma}$（$\gamma$ 为 $M$ 处法线与 $z$ 轴夹角），而 $\cos\gamma=\frac{1}{\sqrt{1+f_x^2+f_y^2}}$，故\textbf{曲面面积元素} $\mathrm{d}A=\sqrt{1+f_x^2+f_y^2}\,\mathrm{d}\sigma$，
\keyformula{A=\iint_D\sqrt{1+\left(\frac{\partial z}{\partial x}\right)^{2}+\left(\frac{\partial z}{\partial y}\right)^{2}}\ \mathrm{d}x\,\mathrm{d}y}
若 $S$：$x=g(y,z)$ 或 $y=h(z,x)$，可分别投影到 $yOz$ 面、$zOx$ 面，把根号内换成相应两偏导的平方和。

\corepoint{公式要求被积函数在 $D$ 上连续。若 $S$ 在投影区域边界处法线趋于与 $z$ 轴垂直（$\cos\gamma\to0$，如球面的"极点"、柱面的母线），被积函数在 $D$ 上无界，不能直接套用，应先取内接区域算出一部分再取极限。}

\subsection{质心}

平面薄片（面密度 $\mu$ 连续）的质量与质心：
\keyformula{M=\iint_D\mu\,\mathrm{d}\sigma,\qquad \bar x=\frac{1}{M}\iint_Dx\mu\,\mathrm{d}\sigma,\qquad \bar y=\frac{1}{M}\iint_Dy\mu\,\mathrm{d}\sigma}
空间物体（体密度 $\mu$ 连续）：$M=\iiint_\Omega\mu\,\mathrm{d}V$，$\bar x=\frac{1}{M}\iiint_\Omega x\mu\,\mathrm{d}V$，$\bar y,\bar z$ 同理。$\mu$ 为常数时质心即\textbf{形心}，$\mu$ 可约去。

\subsection{转动惯量}

平面薄片对 $x$ 轴、$y$ 轴、原点的转动惯量：
\keyformula{I_x=\iint_Dy^{2}\mu\,\mathrm{d}\sigma,\qquad I_y=\iint_Dx^{2}\mu\,\mathrm{d}\sigma,\qquad I_O=I_x+I_y=\iint_D(x^{2}+y^{2})\mu\,\mathrm{d}\sigma}
空间物体：
\[
I_x=\iiint_\Omega(y^2+z^2)\mu\,\mathrm{d}V,\quad
I_y=\iiint_\Omega(z^2+x^2)\mu\,\mathrm{d}V,\quad
I_z=\iiint_\Omega(x^2+y^2)\mu\,\mathrm{d}V,\quad
I_O=I_x+I_y+I_z .
\]
记忆要点：\textbf{对哪根轴取转动惯量，微元里的"距离平方"就是另外两个坐标的平方和}。

\subsection{引力}

设空间物体占有 $\Omega$，体密度 $\mu$ 连续，则它对位于 $M_0(x_0,y_0,z_0)$、质量为 $m$ 的质点的引力 $\boldsymbol F=(F_x,F_y,F_z)$ 为
\[
F_x=km\iiint_\Omega\frac{\mu(x,y,z)(x-x_0)}{r^{3}}\,\mathrm{d}V ,
\]
$F_y,F_z$ 只需把分子换成 $y-y_0$、$z-z_0$，其中 $r=\sqrt{(x-x_0)^2+(y-y_0)^2+(z-z_0)^2}$，$k$ 为引力常数。

\corepoint{分母是 $r^{\mathbf{3}}$ 而非 $r^2$：微元对单位质量的引力大小为 $\frac{k\mu\,\mathrm{d}V}{r^{2}}$，再乘方向余弦 $\frac{x-x_0}{r}$，合起来即 $\frac{x-x_0}{r^3}$。平面薄片对质点的引力把三重积分换成二重积分、$\mathrm{d}V$ 换成 $\mathrm{d}\sigma$，$r$ 仍按空间距离算。}

\section{含参变量的积分}

设 $f(x,y)$ 在矩形 $R=[a,b]\times[c,d]$ 上连续，则 $\Phi(y)=\int_a^bf(x,y)\,\mathrm{d}x$ 确定了 $y$ 在 $[c,d]$ 上的一个函数，称为\textbf{含参变量的积分}（$y$ 为参变量）。

\textbf{定理 1（连续性）}\quad 若 $f(x,y)$ 在 $R$ 上连续，则 $\Phi(y)$ 在 $[c,d]$ 上连续，即 $\lim\limits_{y\to y_0}\int_a^bf(x,y)\,\mathrm{d}x=\int_a^bf(x,y_0)\,\mathrm{d}x$（极限与积分可交换）。

\textbf{定理 2（积分次序交换）}\quad 若 $f$ 在 $R$ 上连续，则
\[
\int_c^d\mathrm{d}y\int_a^bf(x,y)\,\mathrm{d}x=\int_a^b\mathrm{d}x\int_c^df(x,y)\,\mathrm{d}y
\]
（即矩形区域上两种次序的二次积分相等）。

\textbf{定理 3（可导性·莱布尼茨公式）}\quad 若 $f$ 及其偏导数 $f_y$ 都在 $R$ 上连续，则 $\Phi(y)$ 在 $[c,d]$ 上可微，且
\keyformula{\Phi'(y)=\frac{\mathrm{d}}{\mathrm{d}y}\int_a^bf(x,y)\,\mathrm{d}x=\int_a^b\frac{\partial f(x,y)}{\partial y}\,\mathrm{d}x}
即\textbf{积分与求导可交换次序}。

\textbf{定理 4（积分限也含参变量）}\quad 若 $f$ 在 $R$ 上连续，$\alpha(y),\beta(y)$ 可微且 $a\le\alpha,\beta\le b$，则 $\Phi(y)=\int_{\alpha(y)}^{\beta(y)}f(x,y)\,\mathrm{d}x$ 可微，且
\[
\Phi'(y)=\int_{\alpha(y)}^{\beta(y)}f_y(x,y)\,\mathrm{d}x+f(\beta(y),y)\beta'(y)-f(\alpha(y),y)\alpha'(y).
\]

\textbf{常用方法}：(1) \textbf{对参变量微分法}——把待求定积分看作某含参积分的特殊值，用莱布尼茨公式求 $\Phi'(y)$，再积分回 $\Phi(y)$ 并代特殊值；(2) \textbf{先积分后求导}——先把 $\Phi(y)$ 算出再求导；(3) \textbf{交换积分次序}——把一元定积分凑成二重积分再换序。

\corepoint{用莱布尼茨公式前必须核对条件：$f$ 与 $f_y$ 都要在\textbf{闭矩形}上连续。积分限含参变量时别漏后两项。用"对参数微分法"求定积分时，引入参数要使 $\Phi(y_0)$ 好算，且全过程参数不得越出连续性所在的矩形。}
 引入，不含 document 环境

\section{二重积分的概念与性质}
\subsection{引例与定义}

\textbf{引例一：曲顶柱体的体积。}底为 $xOy$ 面上的闭区域 $D$，顶为曲面 $z=f(x,y)$（$f\ge0$ 连续），高随点而变，"体积 $=$ 高 $\times$ 底面积"失效。用曲线网把 $D$ 分成 $n$ 个小闭区域 $\Delta\sigma_i$（既表示小块，也表示其面积），细曲顶柱体近似为以 $f(\xi_i,\eta_i)$ 为高的平顶柱体，求和后令最大直径 $\lambda\to0$ 取极限：$V=\lim\limits_{\lambda\to0}\sum\limits_{i=1}^{n}f(\xi_i,\eta_i)\Delta\sigma_i$。

\textbf{引例二：平面薄片的质量。}薄片占有 $D$，面密度 $\mu(x,y)>0$ 连续，同样"分割—近似—求和—取极限"，得 $m=\lim\limits_{\lambda\to0}\sum\limits_{i=1}^{n}\mu(\xi_i,\eta_i)\Delta\sigma_i$。

两问题的实际意义不同，但所求量都归结为同一形式的和的极限，抽去具体意义即得

\textbf{定义（二重积分）}\quad 设 $f(x,y)$ 是有界闭区域 $D$ 上的有界函数。将 $D$ 任意分成 $n$ 个小闭区域 $\Delta\sigma_i$，在每个 $\Delta\sigma_i$ 上任取一点 $(\xi_i,\eta_i)$，作和 $\sum\limits_{i=1}^{n}f(\xi_i,\eta_i)\Delta\sigma_i$。若当 $\lambda\to0$ 时该和的极限总存在，且与 $D$ 的分法及点 $(\xi_i,\eta_i)$ 的取法无关，则称此极限为 $f(x,y)$ 在 $D$ 上的\textbf{二重积分}，记作
\[
\iint_D f(x,y)\,\mathrm{d}\sigma=\lim_{\lambda\to0}\sum_{i=1}^{n}f(\xi_i,\eta_i)\,\Delta\sigma_i ,
\]
其中 $f(x,y)$ 为\textbf{被积函数}，$f(x,y)\mathrm{d}\sigma$ 为\textbf{被积表达式}，$\mathrm{d}\sigma$ 为\textbf{面积元素}，$x,y$ 为\textbf{积分变量}，$D$ 为\textbf{积分区域}。

\begin{itemize}
  \item 直角坐标中用平行于坐标轴的直线网划分 $D$，除含边界点的小块外都是矩形，$\Delta\sigma=\Delta x\,\Delta y$，故 $\mathrm{d}\sigma=\mathrm{d}x\,\mathrm{d}y$，二重积分也记作 $\iint_Df(x,y)\,\mathrm{d}x\,\mathrm{d}y$。
  \item \textbf{可积性}：$f(x,y)$ 在 $D$ 上\textbf{连续}时积分和极限必存在，即必可积（可积的必要条件是被积函数有界）。教材以后总假定可积。
  \item \textbf{几何意义}：$f\ge0$ 时积分值是以 $D$ 为底、$z=f(x,y)$ 为顶的曲顶柱体体积；$f$ 变号时，积分值等于 $xOy$ 面上方柱体体积减去下方柱体体积之差。
\end{itemize}

\subsection{二重积分的性质}

设所涉积分都存在，$\sigma$ 为 $D$ 的面积。

\begin{enumerate}
  \item \textbf{线性性质}\quad $\iint_D[\alpha f+\beta g]\,\mathrm{d}\sigma=\alpha\iint_Df\,\mathrm{d}\sigma+\beta\iint_Dg\,\mathrm{d}\sigma$（$\alpha,\beta$ 为常数）。
  \item \textbf{对积分区域的可加性}\quad 若 $D$ 被有限条曲线分为有限个部分闭区域（如 $D=D_1\cup D_2$，无公共内点），则 $\iint_Df\,\mathrm{d}\sigma=\iint_{D_1}f\,\mathrm{d}\sigma+\iint_{D_2}f\,\mathrm{d}\sigma$。
  \item \textbf{面积性质}\quad 若在 $D$ 上 $f\equiv1$，则 $\iint_D1\,\mathrm{d}\sigma=\iint_D\mathrm{d}\sigma=\sigma$（高为 $1$ 的平顶柱体体积在数值上等于底面积）。
  \item \textbf{比较性质}\quad 若在 $D$ 上 $f\le g$，则 $\iint_Df\,\mathrm{d}\sigma\le\iint_Dg\,\mathrm{d}\sigma$；特别地 $\left|\iint_Df\,\mathrm{d}\sigma\right|\le\iint_D|f|\,\mathrm{d}\sigma$。
  \item \textbf{估值性质}\quad 设 $M,m$ 为 $f$ 在 $D$ 上的最大值、最小值，则 $m\sigma\le\iint_Df\,\mathrm{d}\sigma\le M\sigma$。
  \item \textbf{中值定理}\quad 设 $f$ 在 $D$ 上连续，则 $D$ 上至少存在一点 $(\xi,\eta)$ 使 $\iint_Df(x,y)\,\mathrm{d}\sigma=f(\xi,\eta)\sigma$。（证：由性质 5 得 $m\le\frac{1}{\sigma}\iint_Df\,\mathrm{d}\sigma\le M$，再用连续函数的介值定理。）
\end{enumerate}

\corepoint{性质 5、6 中的"$\sigma$ 是 $D$ 的面积"不能漏：直角坐标下 $\sigma=\iint_D\mathrm{d}x\mathrm{d}y$，极坐标下 $\sigma=\iint_D\rho\,\mathrm{d}\rho\mathrm{d}\theta$，换元后面积要重算，不能把 $\Delta x\Delta y$ 当面积。另外"$f\le g\Rightarrow$ 积分值 $\le$"只对\textbf{同一积分区域}成立。}

\subsection{对称性简化（常用）}

\begin{itemize}
  \item $D$ 关于 $y$ 轴对称时：
  \[
  \iint_Df(x,y)\,\mathrm{d}\sigma=
  \begin{cases}
  0, & f(-x,y)=-f(x,y),\\[2pt]
  2\displaystyle\iint_{D_1}f(x,y)\,\mathrm{d}\sigma, & f(-x,y)=f(x,y),
  \end{cases}
  \]
  $D_1$ 为 $D$ 中 $x\ge0$ 的部分。
  \item $D$ 关于 $x$ 轴对称时把 $x,y$ 互换；若 $D$ 关于两轴均对称且 $f$ 关于 $x,y$ 都是偶函数，则 $\iint_Df\,\mathrm{d}\sigma=4\iint_{D_1}f\,\mathrm{d}\sigma$（$D_1$ 为第一象限部分）；$D$ 关于原点对称且 $f(-x,-y)=f(x,y)$ 时 $\iint_Df\,\mathrm{d}\sigma=2\iint_{D_{y\ge0}}f\,\mathrm{d}\sigma$。
\end{itemize}

\corepoint{用对称性有\textbf{两个前提}：(1) 积分\textbf{区域}对称；(2) 被积函数在该对称下有确定的奇偶性，缺一不可。判断奇偶性时只有一个变量变号，其余当常数。若 $D$ 关于 $y$ 轴对称而 $f$ 关于 $x$ 为奇函数，则积分为 $0$，这一句常能省掉整段计算。}

\section{二重积分的计算法}

按定义计算只对少数特别简单的被积函数与区域可行。基本方法是把二重积分化为\textbf{两次定积分}（累次积分）。

\subsection{利用直角坐标计算}

取 $\mathrm{d}\sigma=\mathrm{d}x\,\mathrm{d}y$。

\textbf{$X$ 型区域（先对 $y$、后对 $x$）}\quad 若 $D$ 可用不等式 $a\le x\le b$，$\varphi_1(x)\le y\le\varphi_2(x)$ 表示（$\varphi_1\le\varphi_2$ 在 $[a,b]$ 上连续），则
\keyformula{\iint_Df(x,y)\,\mathrm{d}x\,\mathrm{d}y=\int_a^b\mathrm{d}x\int_{\varphi_1(x)}^{\varphi_2(x)}f(x,y)\,\mathrm{d}y}
特征：穿过 $D$ 内部且\textbf{平行于 $y$ 轴}的直线与 $D$ 的边界相交不多于两点。

\textbf{$Y$ 型区域（先对 $x$、后对 $y$）}\quad 若 $D$ 可用 $c\le y\le d$，$\psi_1(y)\le x\le\psi_2(y)$ 表示，则
\keyformula{\iint_Df(x,y)\,\mathrm{d}x\,\mathrm{d}y=\int_c^d\mathrm{d}y\int_{\psi_1(y)}^{\psi_2(y)}f(x,y)\,\mathrm{d}x}
特征：穿过 $D$ 内部且\textbf{平行于 $x$ 轴}的直线与边界相交不多于两点。

两种情形要处理：

\begin{itemize}
  \item 若既有平行于 $y$ 轴的直线与边界相交多于两点、又有平行于 $x$ 轴的直线相交多于两点，则 $D$ 既非 $X$ 型也非 $Y$ 型，须把 $D$ 分成几块，每块为 $X$ 型或 $Y$ 型，再由性质 2 相加。
  \item 若 $D$ 既是 $X$ 型又是 $Y$ 型，则两种次序的二次积分相等，都等于同一个二重积分（这正是累次积分可交换次序的依据）。
\end{itemize}

\textbf{化累次积分的步骤}

\begin{enumerate}
  \item \textbf{画区域}：画出 $D$，求出边界曲线的全部交点。
  \item \textbf{定类型}：判断 $X$ 型、$Y$ 型还是需分块。
  \item \textbf{定限}：外层积分限必为\textbf{常数}，内层积分限是外层变量的函数。
  \item \textbf{算内层}：把外层变量当常数，得到外层变量的函数。
  \item \textbf{算外层}：对剩下的变量作定积分。
\end{enumerate}

\corepoint{定限是唯一的关键，口诀"先画图、后定限，外层常数、内层函数"。常见错误：外层限写成函数；上下限写反；分块时漏块；先对 $y$ 积时把含 $x$ 的因子留在内层（应作为常数提出）。}

\corepoint{\textbf{选择次序}要同时看区域形状与被积函数：既要 $D$ 少分块，又要内层能积出来。含 $\mathrm{e}^{y^2}$、$\frac{\sin y}{y}$、$\sqrt{1+y^3}$ 之类的积分对 $y$ 积不出来，必须\textbf{交换次序}：先由原累次积分还原 $D$ 的不等式组，画图，再按新次序重新定限。这常常是唯一出路。}

\subsection{利用极坐标计算}

当 $D$ 的边界用极坐标表示方便（圆、圆环、扇形）且被积函数含 $x^2+y^2$、$\frac{y}{x}$ 时用极坐标。取 $x=\rho\cos\theta$，$y=\rho\sin\theta$，两族曲线围成的小块面积 $\approx\rho\,\mathrm{d}\rho\,\mathrm{d}\theta$，故\textbf{极坐标下的面积元素} $\mathrm{d}\sigma=\rho\,\mathrm{d}\rho\,\mathrm{d}\theta$，
\keyformula{\iint_Df(x,y)\,\mathrm{d}x\,\mathrm{d}y=\iint_Df(\rho\cos\theta,\rho\sin\theta)\,\rho\,\mathrm{d}\rho\,\mathrm{d}\theta}
即：$x,y$ 分别换成 $\rho\cos\theta,\rho\sin\theta$，$\mathrm{d}x\mathrm{d}y$ 换成 $\rho\,\mathrm{d}\rho\mathrm{d}\theta$（\textbf{这个 $\rho$ 不能漏}）。

设 $D$ 为 $\rho_1(\theta)\le\rho\le\rho_2(\theta)$，$\alpha\le\theta\le\beta$，则
\keyformula{\iint_Df(\rho\cos\theta,\rho\sin\theta)\,\rho\,\mathrm{d}\rho\,\mathrm{d}\theta=\int_\alpha^\beta\mathrm{d}\theta\int_{\rho_1(\theta)}^{\rho_2(\theta)}f(\rho\cos\theta,\rho\sin\theta)\,\rho\,\mathrm{d}\rho}
两种特例：极点在 $D$ 的\textbf{边界}上（曲边扇形）时 $\rho_1\equiv0$，$\rho$ 的上限为 $\rho(\theta)$；极点在 $D$ 的\textbf{内部}时 $0\le\rho\le\rho(\theta)$，$0\le\theta\le2\pi$。由性质 3 得曲边扇形的面积 $\sigma=\frac{1}{2}\int_\alpha^\beta[\rho_2^2(\theta)-\rho_1^2(\theta)]\,\mathrm{d}\theta$。

\corepoint{极坐标三个高频错误：(1) 漏掉面积元素里的 $\rho$；(2) 不管区域张多大角，一律写 $[0,2\pi]$；(3) 内外层限的位置搞反。另需记准两个常用转换：$x^2+y^2=2ax$ 即 $\rho=2a\cos\theta$（$\theta\in[-\frac{\pi}{2},\frac{\pi}{2}]$），$x^2+y^2=2ay$ 即 $\rho=2a\sin\theta$（$\theta\in[0,\pi]$）。}

\subsection{二重积分的换元法}

设 $f(x,y)$ 在 $D$ 上连续，变换 $x=x(u,v)$，$y=y(u,v)$ 把 $uOv$ 面上的闭区域 $D'$ 一一对应地映成 $D$，$x,y$ 在 $D'$ 上有一阶连续偏导数，其\textbf{雅可比行列式}
\[
J(u,v)=\frac{\partial(x,y)}{\partial(u,v)}
=\begin{vmatrix}
\dfrac{\partial x}{\partial u} & \dfrac{\partial x}{\partial v}\\[8pt]
\dfrac{\partial y}{\partial u} & \dfrac{\partial y}{\partial v}
\end{vmatrix}\ne0 ,
\]
则
\keyformula{\iint_Df(x,y)\,\mathrm{d}x\,\mathrm{d}y=\iint_{D'}f\bigl(x(u,v),y(u,v)\bigr)\,\bigl|J(u,v)\bigr|\,\mathrm{d}u\,\mathrm{d}v}
极坐标 $x=\rho\cos\theta,y=\rho\sin\theta$ 是其特例，$|J|=\rho$。常用\textbf{广义极坐标}（椭圆域 $\frac{x^2}{a^2}+\frac{y^2}{b^2}\le1$）：$x=a\rho\cos\theta$，$y=b\rho\sin\theta$，$|J|=ab\rho$。

\corepoint{公式里是 $|J|$ 的\textbf{绝对值}；换元后区域一并换成 $D'$，积分限须按 $D'$ 重定；变换须一一对应（$J\ne0$ 只保证局部一一），否则分块。}

\section{三重积分}
\subsection{概念}

\textbf{定义}\quad 设 $f(x,y,z)$ 是空间有界闭区域 $\Omega$ 上的有界函数。把 $\Omega$ 任意分成 $n$ 个小闭区域 $\Delta V_i$（也表示其体积），任取 $(\xi_i,\eta_i,\zeta_i)$ 作和 $\sum f(\xi_i,\eta_i,\zeta_i)\Delta V_i$。若 $\lambda\to0$（各小块直径最大值）时极限总存在且与分法、取法无关，则称之为 $f$ 在 $\Omega$ 上的\textbf{三重积分}：
\[
\iiint_\Omega f(x,y,z)\,\mathrm{d}V=\lim_{\lambda\to0}\sum_{i=1}^{n}f(\xi_i,\eta_i,\zeta_i)\,\Delta V_i ,
\]
$\mathrm{d}V$ 为\textbf{体积元素}；直角坐标下 $\mathrm{d}V=\mathrm{d}x\,\mathrm{d}y\,\mathrm{d}z$。$f$ 连续时必可积；$f\ge0$ 时可解释为空间物体的质量（$f$ 为体密度）。

性质与二重积分完全平行：线性、对区域可加、$\iiint_\Omega\mathrm{d}V=V$、比较、估值、中值定理 $\iiint_\Omega f\,\mathrm{d}V=f(\xi,\eta,\zeta)V$。

\subsection{直角坐标：投影法与截面法}

\textbf{投影法（先一后二）}\quad 设平行于 $z$ 轴且穿过 $\Omega$ 内部的直线与边界曲面 $S$ 相交不多于两点。把 $\Omega$ 投影到 $xOy$ 面得 $D_{xy}$，$S$ 被分成下、上两部分 $S_1:z=z_1(x,y)$，$S_2:z=z_2(x,y)$（$z_1\le z_2$ 均在 $D_{xy}$ 上连续），则
\keyformula{\iiint_\Omega f\,\mathrm{d}V=\iint_{D_{xy}}\mathrm{d}x\,\mathrm{d}y\int_{z_1(x,y)}^{z_2(x,y)}f(x,y,z)\,\mathrm{d}z}
即先对 $z$、再对 $x,y$ 的三次积分。若平行于 $x$ 轴或 $y$ 轴的直线与 $S$ 相交不多于两点，也可投影到 $yOz$ 面或 $zOx$ 面，化为其他顺序的三次积分；交点多于两点时把 $\Omega$ 分块相加。

\textbf{截面法（先二后一）}\quad 设 $\Omega=\{(x,y,z)\mid c\le z\le d,\ (x,y)\in D_z\}$，其中 $D_z$ 是用竖坐标 $z$ 的平面截 $\Omega$ 所得的平面闭区域，则
\keyformula{\iiint_\Omega f\,\mathrm{d}V=\int_c^d\mathrm{d}z\iint_{D_z}f(x,y,z)\,\mathrm{d}x\,\mathrm{d}y}
当被积函数只含 $z$（或含 $z$ 的因子可分离）且截面面积易求时，内层二重积分 $=$ 截面面积 $\times$ 该因子。

\corepoint{两法的差别只在"先算哪个"：投影法是\textbf{先一后二}，要求沿 $z$ 方向穿入穿出各一次；截面法是\textbf{先二后一}，要求截面形状随 $z$ 变化简单。判据：被积函数中 $z$ 的因子能否提出内层——能提就优先截面法。}

\subsection{柱面坐标}

设点 $M$ 在 $xOy$ 面上投影 $P$ 的极坐标为 $(\rho,\theta)$，则 $(\rho,\theta,z)$ 为 $M$ 的\textbf{柱面坐标}，$0\le\rho<+\infty$，$0\le\theta\le2\pi$，$-\infty<z<+\infty$。三组坐标面：$\rho=$ 常数（以 $z$ 轴为轴的圆柱面）、$\theta=$ 常数（过 $z$ 轴的半平面）、$z=$ 常数（平行于 $xOy$ 面的平面）。关系式与体积元素：
\keyformula{x=\rho\cos\theta,\quad y=\rho\sin\theta,\quad z=z;\qquad \mathrm{d}V=\rho\,\mathrm{d}\rho\,\mathrm{d}\theta\,\mathrm{d}z}
\keyformula{\iiint_\Omega f\,\mathrm{d}V=\iiint_{\Omega'}f(\rho\cos\theta,\rho\sin\theta,z)\,\rho\,\mathrm{d}\rho\,\mathrm{d}\theta\,\mathrm{d}z}
再化为三次积分，积分限按 $\rho,\theta,z$ 在 $\Omega$ 中的变化范围确定。

\corepoint{柱面坐标适合"绕 $z$ 轴旋转"的区域（圆柱、圆锥、旋转抛物面、圆环）与被积函数含 $x^2+y^2$ 的情形。$\rho$ 的限与二重积分的极坐标一致：先看投影区域张多大角，再看沿射线 $\rho$ 从哪进哪出；别忘了因子 $\rho$。}

\subsection{球面坐标}

设 $r$ 为原点 $O$ 与 $M$ 的距离，$\varphi$ 为有向线段 $\overrightarrow{OM}$ 与 $z$ 轴正向的夹角，$\theta$ 为从 $z$ 轴正向来看自 $x$ 轴按逆时针转到 $\overrightarrow{OP}$ 的角（$P$ 为 $M$ 在 $xOy$ 面上的投影），则 $(r,\varphi,\theta)$ 为 $M$ 的\textbf{球面坐标}，$0\le r<+\infty$，$0\le\varphi\le\pi$，$0\le\theta\le2\pi$。三组坐标面：$r=$ 常数（以原点为球心的球面）、$\varphi=$ 常数（以原点为顶点、$z$ 轴为轴的圆锥面）、$\theta=$ 常数（过 $z$ 轴的半平面）。
\keyformula{x=r\sin\varphi\cos\theta,\quad y=r\sin\varphi\sin\theta,\quad z=r\cos\varphi}
\keyformula{\mathrm{d}V=r^{2}\sin\varphi\,\mathrm{d}r\,\mathrm{d}\varphi\,\mathrm{d}\theta}
\keyformula{\iiint_\Omega f\,\mathrm{d}V=\iiint_{\Omega'}f(r\sin\varphi\cos\theta,\ r\sin\varphi\sin\theta,\ r\cos\varphi)\,r^{2}\sin\varphi\,\mathrm{d}r\,\mathrm{d}\varphi\,\mathrm{d}\theta}
一般化为先对 $r$、再对 $\varphi$、最后对 $\theta$ 的三次积分：$0\le\theta\le2\pi$，$0\le\varphi\le\pi$（只在上半球时 $0\le\varphi\le\frac{\pi}{2}$），$r$ 的限由"自原点出发的射线穿过 $\Omega$"确定。

\corepoint{两个最易错处：(1) 体积元素是 $r^2\sin\varphi$，不是 $\cos\varphi$、也不是 $\sin\theta$；(2) $\varphi$ 是\textbf{与 $z$ 轴正向的夹角}，不是与 $xOy$ 面的夹角——上半球对应 $\varphi\in[0,\frac{\pi}{2}]$，圆锥面 $z=\sqrt{x^2+y^2}$ 对应 $\varphi=\frac{\pi}{4}$。记不住就用 $z=r\cos\varphi$ 现推。}

\subsection{三种坐标的选择}

\begin{center}
\begin{tabular}{@{}lll@{}}
\toprule
坐标 & 适用的积分区域 & 体积元素 \\
\midrule
直角坐标 & 长方体、四面体、由平面与柱面围成的可分层区域 & $\mathrm{d}x\,\mathrm{d}y\,\mathrm{d}z$ \\
柱面坐标 & 圆柱、圆锥、旋转抛物面、绕 $z$ 轴的旋转体 & $\rho\,\mathrm{d}\rho\,\mathrm{d}\theta\,\mathrm{d}z$ \\
球面坐标 & 球体、球壳、球与锥面所围的"冰淇淋"形 & $r^{2}\sin\varphi\,\mathrm{d}r\,\mathrm{d}\varphi\,\mathrm{d}\theta$ \\
\bottomrule
\end{tabular}
\end{center}

选择原则：先看\textbf{边界曲面}（球用球面坐标、柱面用柱面坐标），再看\textbf{被积函数}含 $x^2+y^2$（柱面坐标）还是 $x^2+y^2+z^2$（球面坐标）；冲突时以"积分限好定"为准。

\section{重积分的应用}
\subsection{元素法}

若所求量 $U$ (1) 对积分区域具有可加性，(2) 在区域内某点的很小邻域上其部分量可用 $f(M)\,\mathrm{d}\Omega$ 近似（$\mathrm{d}\Omega$ 为面积或体积元素），则 $U=\iint_Df\,\mathrm{d}\sigma$ 或 $U=\iiint_\Omega f\,\mathrm{d}V$。关键是写出\textbf{正确的微元}：几何量用面积（体积）元素，物理量用"密度 $\times$ 元素"。

\subsection{曲面的面积}

设曲面 $S:z=f(x,y)$，$f$ 在 $D$（$S$ 在 $xOy$ 面上的投影区域）有连续偏导。$S$ 上小片面积 $\mathrm{d}A$ 与其投影 $\mathrm{d}\sigma$ 之比为 $\frac{1}{\cos\gamma}$（$\gamma$ 为 $M$ 处法线与 $z$ 轴夹角），而 $\cos\gamma=\frac{1}{\sqrt{1+f_x^2+f_y^2}}$，故\textbf{曲面面积元素} $\mathrm{d}A=\sqrt{1+f_x^2+f_y^2}\,\mathrm{d}\sigma$，
\keyformula{A=\iint_D\sqrt{1+\left(\frac{\partial z}{\partial x}\right)^{2}+\left(\frac{\partial z}{\partial y}\right)^{2}}\ \mathrm{d}x\,\mathrm{d}y}
若 $S$：$x=g(y,z)$ 或 $y=h(z,x)$，可分别投影到 $yOz$ 面、$zOx$ 面，把根号内换成相应两偏导的平方和。

\corepoint{公式要求被积函数在 $D$ 上连续。若 $S$ 在投影区域边界处法线趋于与 $z$ 轴垂直（$\cos\gamma\to0$，如球面的"极点"、柱面的母线），被积函数在 $D$ 上无界，不能直接套用，应先取内接区域算出一部分再取极限。}

\subsection{质心}

平面薄片（面密度 $\mu$ 连续）的质量与质心：
\keyformula{M=\iint_D\mu\,\mathrm{d}\sigma,\qquad \bar x=\frac{1}{M}\iint_Dx\mu\,\mathrm{d}\sigma,\qquad \bar y=\frac{1}{M}\iint_Dy\mu\,\mathrm{d}\sigma}
空间物体（体密度 $\mu$ 连续）：$M=\iiint_\Omega\mu\,\mathrm{d}V$，$\bar x=\frac{1}{M}\iiint_\Omega x\mu\,\mathrm{d}V$，$\bar y,\bar z$ 同理。$\mu$ 为常数时质心即\textbf{形心}，$\mu$ 可约去。

\subsection{转动惯量}

平面薄片对 $x$ 轴、$y$ 轴、原点的转动惯量：
\keyformula{I_x=\iint_Dy^{2}\mu\,\mathrm{d}\sigma,\qquad I_y=\iint_Dx^{2}\mu\,\mathrm{d}\sigma,\qquad I_O=I_x+I_y=\iint_D(x^{2}+y^{2})\mu\,\mathrm{d}\sigma}
空间物体：
\[
I_x=\iiint_\Omega(y^2+z^2)\mu\,\mathrm{d}V,\quad
I_y=\iiint_\Omega(z^2+x^2)\mu\,\mathrm{d}V,\quad
I_z=\iiint_\Omega(x^2+y^2)\mu\,\mathrm{d}V,\quad
I_O=I_x+I_y+I_z .
\]
记忆要点：\textbf{对哪根轴取转动惯量，微元里的"距离平方"就是另外两个坐标的平方和}。

\subsection{引力}

设空间物体占有 $\Omega$，体密度 $\mu$ 连续，则它对位于 $M_0(x_0,y_0,z_0)$、质量为 $m$ 的质点的引力 $\boldsymbol F=(F_x,F_y,F_z)$ 为
\[
F_x=km\iiint_\Omega\frac{\mu(x,y,z)(x-x_0)}{r^{3}}\,\mathrm{d}V ,
\]
$F_y,F_z$ 只需把分子换成 $y-y_0$、$z-z_0$，其中 $r=\sqrt{(x-x_0)^2+(y-y_0)^2+(z-z_0)^2}$，$k$ 为引力常数。

\corepoint{分母是 $r^{\mathbf{3}}$ 而非 $r^2$：微元对单位质量的引力大小为 $\frac{k\mu\,\mathrm{d}V}{r^{2}}$，再乘方向余弦 $\frac{x-x_0}{r}$，合起来即 $\frac{x-x_0}{r^3}$。平面薄片对质点的引力把三重积分换成二重积分、$\mathrm{d}V$ 换成 $\mathrm{d}\sigma$，$r$ 仍按空间距离算。}

\section{含参变量的积分}

设 $f(x,y)$ 在矩形 $R=[a,b]\times[c,d]$ 上连续，则 $\Phi(y)=\int_a^bf(x,y)\,\mathrm{d}x$ 确定了 $y$ 在 $[c,d]$ 上的一个函数，称为\textbf{含参变量的积分}（$y$ 为参变量）。

\textbf{定理 1（连续性）}\quad 若 $f(x,y)$ 在 $R$ 上连续，则 $\Phi(y)$ 在 $[c,d]$ 上连续，即 $\lim\limits_{y\to y_0}\int_a^bf(x,y)\,\mathrm{d}x=\int_a^bf(x,y_0)\,\mathrm{d}x$（极限与积分可交换）。

\textbf{定理 2（积分次序交换）}\quad 若 $f$ 在 $R$ 上连续，则
\[
\int_c^d\mathrm{d}y\int_a^bf(x,y)\,\mathrm{d}x=\int_a^b\mathrm{d}x\int_c^df(x,y)\,\mathrm{d}y
\]
（即矩形区域上两种次序的二次积分相等）。

\textbf{定理 3（可导性·莱布尼茨公式）}\quad 若 $f$ 及其偏导数 $f_y$ 都在 $R$ 上连续，则 $\Phi(y)$ 在 $[c,d]$ 上可微，且
\keyformula{\Phi'(y)=\frac{\mathrm{d}}{\mathrm{d}y}\int_a^bf(x,y)\,\mathrm{d}x=\int_a^b\frac{\partial f(x,y)}{\partial y}\,\mathrm{d}x}
即\textbf{积分与求导可交换次序}。

\textbf{定理 4（积分限也含参变量）}\quad 若 $f$ 在 $R$ 上连续，$\alpha(y),\beta(y)$ 可微且 $a\le\alpha,\beta\le b$，则 $\Phi(y)=\int_{\alpha(y)}^{\beta(y)}f(x,y)\,\mathrm{d}x$ 可微，且
\[
\Phi'(y)=\int_{\alpha(y)}^{\beta(y)}f_y(x,y)\,\mathrm{d}x+f(\beta(y),y)\beta'(y)-f(\alpha(y),y)\alpha'(y).
\]

\textbf{常用方法}：(1) \textbf{对参变量微分法}——把待求定积分看作某含参积分的特殊值，用莱布尼茨公式求 $\Phi'(y)$，再积分回 $\Phi(y)$ 并代特殊值；(2) \textbf{先积分后求导}——先把 $\Phi(y)$ 算出再求导；(3) \textbf{交换积分次序}——把一元定积分凑成二重积分再换序。

\corepoint{用莱布尼茨公式前必须核对条件：$f$ 与 $f_y$ 都要在\textbf{闭矩形}上连续。积分限含参变量时别漏后两项。用"对参数微分法"求定积分时，引入参数要使 $\Phi(y_0)$ 好算，且全过程参数不得越出连续性所在的矩形。}
 引入，不含 document 环境

\section{二重积分的概念与性质}
\subsection{引例与定义}

\textbf{引例一：曲顶柱体的体积。}底为 $xOy$ 面上的闭区域 $D$，顶为曲面 $z=f(x,y)$（$f\ge0$ 连续），高随点而变，"体积 $=$ 高 $\times$ 底面积"失效。用曲线网把 $D$ 分成 $n$ 个小闭区域 $\Delta\sigma_i$（既表示小块，也表示其面积），细曲顶柱体近似为以 $f(\xi_i,\eta_i)$ 为高的平顶柱体，求和后令最大直径 $\lambda\to0$ 取极限：$V=\lim\limits_{\lambda\to0}\sum\limits_{i=1}^{n}f(\xi_i,\eta_i)\Delta\sigma_i$。

\textbf{引例二：平面薄片的质量。}薄片占有 $D$，面密度 $\mu(x,y)>0$ 连续，同样"分割—近似—求和—取极限"，得 $m=\lim\limits_{\lambda\to0}\sum\limits_{i=1}^{n}\mu(\xi_i,\eta_i)\Delta\sigma_i$。

两问题的实际意义不同，但所求量都归结为同一形式的和的极限，抽去具体意义即得

\textbf{定义（二重积分）}\quad 设 $f(x,y)$ 是有界闭区域 $D$ 上的有界函数。将 $D$ 任意分成 $n$ 个小闭区域 $\Delta\sigma_i$，在每个 $\Delta\sigma_i$ 上任取一点 $(\xi_i,\eta_i)$，作和 $\sum\limits_{i=1}^{n}f(\xi_i,\eta_i)\Delta\sigma_i$。若当 $\lambda\to0$ 时该和的极限总存在，且与 $D$ 的分法及点 $(\xi_i,\eta_i)$ 的取法无关，则称此极限为 $f(x,y)$ 在 $D$ 上的\textbf{二重积分}，记作
\[
\iint_D f(x,y)\,\mathrm{d}\sigma=\lim_{\lambda\to0}\sum_{i=1}^{n}f(\xi_i,\eta_i)\,\Delta\sigma_i ,
\]
其中 $f(x,y)$ 为\textbf{被积函数}，$f(x,y)\mathrm{d}\sigma$ 为\textbf{被积表达式}，$\mathrm{d}\sigma$ 为\textbf{面积元素}，$x,y$ 为\textbf{积分变量}，$D$ 为\textbf{积分区域}。

\begin{itemize}
  \item 直角坐标中用平行于坐标轴的直线网划分 $D$，除含边界点的小块外都是矩形，$\Delta\sigma=\Delta x\,\Delta y$，故 $\mathrm{d}\sigma=\mathrm{d}x\,\mathrm{d}y$，二重积分也记作 $\iint_Df(x,y)\,\mathrm{d}x\,\mathrm{d}y$。
  \item \textbf{可积性}：$f(x,y)$ 在 $D$ 上\textbf{连续}时积分和极限必存在，即必可积（可积的必要条件是被积函数有界）。教材以后总假定可积。
  \item \textbf{几何意义}：$f\ge0$ 时积分值是以 $D$ 为底、$z=f(x,y)$ 为顶的曲顶柱体体积；$f$ 变号时，积分值等于 $xOy$ 面上方柱体体积减去下方柱体体积之差。
\end{itemize}

\subsection{二重积分的性质}

设所涉积分都存在，$\sigma$ 为 $D$ 的面积。

\begin{enumerate}
  \item \textbf{线性性质}\quad $\iint_D[\alpha f+\beta g]\,\mathrm{d}\sigma=\alpha\iint_Df\,\mathrm{d}\sigma+\beta\iint_Dg\,\mathrm{d}\sigma$（$\alpha,\beta$ 为常数）。
  \item \textbf{对积分区域的可加性}\quad 若 $D$ 被有限条曲线分为有限个部分闭区域（如 $D=D_1\cup D_2$，无公共内点），则 $\iint_Df\,\mathrm{d}\sigma=\iint_{D_1}f\,\mathrm{d}\sigma+\iint_{D_2}f\,\mathrm{d}\sigma$。
  \item \textbf{面积性质}\quad 若在 $D$ 上 $f\equiv1$，则 $\iint_D1\,\mathrm{d}\sigma=\iint_D\mathrm{d}\sigma=\sigma$（高为 $1$ 的平顶柱体体积在数值上等于底面积）。
  \item \textbf{比较性质}\quad 若在 $D$ 上 $f\le g$，则 $\iint_Df\,\mathrm{d}\sigma\le\iint_Dg\,\mathrm{d}\sigma$；特别地 $\left|\iint_Df\,\mathrm{d}\sigma\right|\le\iint_D|f|\,\mathrm{d}\sigma$。
  \item \textbf{估值性质}\quad 设 $M,m$ 为 $f$ 在 $D$ 上的最大值、最小值，则 $m\sigma\le\iint_Df\,\mathrm{d}\sigma\le M\sigma$。
  \item \textbf{中值定理}\quad 设 $f$ 在 $D$ 上连续，则 $D$ 上至少存在一点 $(\xi,\eta)$ 使 $\iint_Df(x,y)\,\mathrm{d}\sigma=f(\xi,\eta)\sigma$。（证：由性质 5 得 $m\le\frac{1}{\sigma}\iint_Df\,\mathrm{d}\sigma\le M$，再用连续函数的介值定理。）
\end{enumerate}

\corepoint{性质 5、6 中的"$\sigma$ 是 $D$ 的面积"不能漏：直角坐标下 $\sigma=\iint_D\mathrm{d}x\mathrm{d}y$，极坐标下 $\sigma=\iint_D\rho\,\mathrm{d}\rho\mathrm{d}\theta$，换元后面积要重算，不能把 $\Delta x\Delta y$ 当面积。另外"$f\le g\Rightarrow$ 积分值 $\le$"只对\textbf{同一积分区域}成立。}

\subsection{对称性简化（常用）}

\begin{itemize}
  \item $D$ 关于 $y$ 轴对称时：
  \[
  \iint_Df(x,y)\,\mathrm{d}\sigma=
  \begin{cases}
  0, & f(-x,y)=-f(x,y),\\[2pt]
  2\displaystyle\iint_{D_1}f(x,y)\,\mathrm{d}\sigma, & f(-x,y)=f(x,y),
  \end{cases}
  \]
  $D_1$ 为 $D$ 中 $x\ge0$ 的部分。
  \item $D$ 关于 $x$ 轴对称时把 $x,y$ 互换；若 $D$ 关于两轴均对称且 $f$ 关于 $x,y$ 都是偶函数，则 $\iint_Df\,\mathrm{d}\sigma=4\iint_{D_1}f\,\mathrm{d}\sigma$（$D_1$ 为第一象限部分）；$D$ 关于原点对称且 $f(-x,-y)=f(x,y)$ 时 $\iint_Df\,\mathrm{d}\sigma=2\iint_{D_{y\ge0}}f\,\mathrm{d}\sigma$。
\end{itemize}

\corepoint{用对称性有\textbf{两个前提}：(1) 积分\textbf{区域}对称；(2) 被积函数在该对称下有确定的奇偶性，缺一不可。判断奇偶性时只有一个变量变号，其余当常数。若 $D$ 关于 $y$ 轴对称而 $f$ 关于 $x$ 为奇函数，则积分为 $0$，这一句常能省掉整段计算。}

\section{二重积分的计算法}

按定义计算只对少数特别简单的被积函数与区域可行。基本方法是把二重积分化为\textbf{两次定积分}（累次积分）。

\subsection{利用直角坐标计算}

取 $\mathrm{d}\sigma=\mathrm{d}x\,\mathrm{d}y$。

\textbf{$X$ 型区域（先对 $y$、后对 $x$）}\quad 若 $D$ 可用不等式 $a\le x\le b$，$\varphi_1(x)\le y\le\varphi_2(x)$ 表示（$\varphi_1\le\varphi_2$ 在 $[a,b]$ 上连续），则
\keyformula{\iint_Df(x,y)\,\mathrm{d}x\,\mathrm{d}y=\int_a^b\mathrm{d}x\int_{\varphi_1(x)}^{\varphi_2(x)}f(x,y)\,\mathrm{d}y}
特征：穿过 $D$ 内部且\textbf{平行于 $y$ 轴}的直线与 $D$ 的边界相交不多于两点。

\textbf{$Y$ 型区域（先对 $x$、后对 $y$）}\quad 若 $D$ 可用 $c\le y\le d$，$\psi_1(y)\le x\le\psi_2(y)$ 表示，则
\keyformula{\iint_Df(x,y)\,\mathrm{d}x\,\mathrm{d}y=\int_c^d\mathrm{d}y\int_{\psi_1(y)}^{\psi_2(y)}f(x,y)\,\mathrm{d}x}
特征：穿过 $D$ 内部且\textbf{平行于 $x$ 轴}的直线与边界相交不多于两点。

两种情形要处理：

\begin{itemize}
  \item 若既有平行于 $y$ 轴的直线与边界相交多于两点、又有平行于 $x$ 轴的直线相交多于两点，则 $D$ 既非 $X$ 型也非 $Y$ 型，须把 $D$ 分成几块，每块为 $X$ 型或 $Y$ 型，再由性质 2 相加。
  \item 若 $D$ 既是 $X$ 型又是 $Y$ 型，则两种次序的二次积分相等，都等于同一个二重积分（这正是累次积分可交换次序的依据）。
\end{itemize}

\textbf{化累次积分的步骤}

\begin{enumerate}
  \item \textbf{画区域}：画出 $D$，求出边界曲线的全部交点。
  \item \textbf{定类型}：判断 $X$ 型、$Y$ 型还是需分块。
  \item \textbf{定限}：外层积分限必为\textbf{常数}，内层积分限是外层变量的函数。
  \item \textbf{算内层}：把外层变量当常数，得到外层变量的函数。
  \item \textbf{算外层}：对剩下的变量作定积分。
\end{enumerate}

\corepoint{定限是唯一的关键，口诀"先画图、后定限，外层常数、内层函数"。常见错误：外层限写成函数；上下限写反；分块时漏块；先对 $y$ 积时把含 $x$ 的因子留在内层（应作为常数提出）。}

\corepoint{\textbf{选择次序}要同时看区域形状与被积函数：既要 $D$ 少分块，又要内层能积出来。含 $\mathrm{e}^{y^2}$、$\frac{\sin y}{y}$、$\sqrt{1+y^3}$ 之类的积分对 $y$ 积不出来，必须\textbf{交换次序}：先由原累次积分还原 $D$ 的不等式组，画图，再按新次序重新定限。这常常是唯一出路。}

\subsection{利用极坐标计算}

当 $D$ 的边界用极坐标表示方便（圆、圆环、扇形）且被积函数含 $x^2+y^2$、$\frac{y}{x}$ 时用极坐标。取 $x=\rho\cos\theta$，$y=\rho\sin\theta$，两族曲线围成的小块面积 $\approx\rho\,\mathrm{d}\rho\,\mathrm{d}\theta$，故\textbf{极坐标下的面积元素} $\mathrm{d}\sigma=\rho\,\mathrm{d}\rho\,\mathrm{d}\theta$，
\keyformula{\iint_Df(x,y)\,\mathrm{d}x\,\mathrm{d}y=\iint_Df(\rho\cos\theta,\rho\sin\theta)\,\rho\,\mathrm{d}\rho\,\mathrm{d}\theta}
即：$x,y$ 分别换成 $\rho\cos\theta,\rho\sin\theta$，$\mathrm{d}x\mathrm{d}y$ 换成 $\rho\,\mathrm{d}\rho\mathrm{d}\theta$（\textbf{这个 $\rho$ 不能漏}）。

设 $D$ 为 $\rho_1(\theta)\le\rho\le\rho_2(\theta)$，$\alpha\le\theta\le\beta$，则
\keyformula{\iint_Df(\rho\cos\theta,\rho\sin\theta)\,\rho\,\mathrm{d}\rho\,\mathrm{d}\theta=\int_\alpha^\beta\mathrm{d}\theta\int_{\rho_1(\theta)}^{\rho_2(\theta)}f(\rho\cos\theta,\rho\sin\theta)\,\rho\,\mathrm{d}\rho}
两种特例：极点在 $D$ 的\textbf{边界}上（曲边扇形）时 $\rho_1\equiv0$，$\rho$ 的上限为 $\rho(\theta)$；极点在 $D$ 的\textbf{内部}时 $0\le\rho\le\rho(\theta)$，$0\le\theta\le2\pi$。由性质 3 得曲边扇形的面积 $\sigma=\frac{1}{2}\int_\alpha^\beta[\rho_2^2(\theta)-\rho_1^2(\theta)]\,\mathrm{d}\theta$。

\corepoint{极坐标三个高频错误：(1) 漏掉面积元素里的 $\rho$；(2) 不管区域张多大角，一律写 $[0,2\pi]$；(3) 内外层限的位置搞反。另需记准两个常用转换：$x^2+y^2=2ax$ 即 $\rho=2a\cos\theta$（$\theta\in[-\frac{\pi}{2},\frac{\pi}{2}]$），$x^2+y^2=2ay$ 即 $\rho=2a\sin\theta$（$\theta\in[0,\pi]$）。}

\subsection{二重积分的换元法}

设 $f(x,y)$ 在 $D$ 上连续，变换 $x=x(u,v)$，$y=y(u,v)$ 把 $uOv$ 面上的闭区域 $D'$ 一一对应地映成 $D$，$x,y$ 在 $D'$ 上有一阶连续偏导数，其\textbf{雅可比行列式}
\[
J(u,v)=\frac{\partial(x,y)}{\partial(u,v)}
=\begin{vmatrix}
\dfrac{\partial x}{\partial u} & \dfrac{\partial x}{\partial v}\\[8pt]
\dfrac{\partial y}{\partial u} & \dfrac{\partial y}{\partial v}
\end{vmatrix}\ne0 ,
\]
则
\keyformula{\iint_Df(x,y)\,\mathrm{d}x\,\mathrm{d}y=\iint_{D'}f\bigl(x(u,v),y(u,v)\bigr)\,\bigl|J(u,v)\bigr|\,\mathrm{d}u\,\mathrm{d}v}
极坐标 $x=\rho\cos\theta,y=\rho\sin\theta$ 是其特例，$|J|=\rho$。常用\textbf{广义极坐标}（椭圆域 $\frac{x^2}{a^2}+\frac{y^2}{b^2}\le1$）：$x=a\rho\cos\theta$，$y=b\rho\sin\theta$，$|J|=ab\rho$。

\corepoint{公式里是 $|J|$ 的\textbf{绝对值}；换元后区域一并换成 $D'$，积分限须按 $D'$ 重定；变换须一一对应（$J\ne0$ 只保证局部一一），否则分块。}

\section{三重积分}
\subsection{概念}

\textbf{定义}\quad 设 $f(x,y,z)$ 是空间有界闭区域 $\Omega$ 上的有界函数。把 $\Omega$ 任意分成 $n$ 个小闭区域 $\Delta V_i$（也表示其体积），任取 $(\xi_i,\eta_i,\zeta_i)$ 作和 $\sum f(\xi_i,\eta_i,\zeta_i)\Delta V_i$。若 $\lambda\to0$（各小块直径最大值）时极限总存在且与分法、取法无关，则称之为 $f$ 在 $\Omega$ 上的\textbf{三重积分}：
\[
\iiint_\Omega f(x,y,z)\,\mathrm{d}V=\lim_{\lambda\to0}\sum_{i=1}^{n}f(\xi_i,\eta_i,\zeta_i)\,\Delta V_i ,
\]
$\mathrm{d}V$ 为\textbf{体积元素}；直角坐标下 $\mathrm{d}V=\mathrm{d}x\,\mathrm{d}y\,\mathrm{d}z$。$f$ 连续时必可积；$f\ge0$ 时可解释为空间物体的质量（$f$ 为体密度）。

性质与二重积分完全平行：线性、对区域可加、$\iiint_\Omega\mathrm{d}V=V$、比较、估值、中值定理 $\iiint_\Omega f\,\mathrm{d}V=f(\xi,\eta,\zeta)V$。

\subsection{直角坐标：投影法与截面法}

\textbf{投影法（先一后二）}\quad 设平行于 $z$ 轴且穿过 $\Omega$ 内部的直线与边界曲面 $S$ 相交不多于两点。把 $\Omega$ 投影到 $xOy$ 面得 $D_{xy}$，$S$ 被分成下、上两部分 $S_1:z=z_1(x,y)$，$S_2:z=z_2(x,y)$（$z_1\le z_2$ 均在 $D_{xy}$ 上连续），则
\keyformula{\iiint_\Omega f\,\mathrm{d}V=\iint_{D_{xy}}\mathrm{d}x\,\mathrm{d}y\int_{z_1(x,y)}^{z_2(x,y)}f(x,y,z)\,\mathrm{d}z}
即先对 $z$、再对 $x,y$ 的三次积分。若平行于 $x$ 轴或 $y$ 轴的直线与 $S$ 相交不多于两点，也可投影到 $yOz$ 面或 $zOx$ 面，化为其他顺序的三次积分；交点多于两点时把 $\Omega$ 分块相加。

\textbf{截面法（先二后一）}\quad 设 $\Omega=\{(x,y,z)\mid c\le z\le d,\ (x,y)\in D_z\}$，其中 $D_z$ 是用竖坐标 $z$ 的平面截 $\Omega$ 所得的平面闭区域，则
\keyformula{\iiint_\Omega f\,\mathrm{d}V=\int_c^d\mathrm{d}z\iint_{D_z}f(x,y,z)\,\mathrm{d}x\,\mathrm{d}y}
当被积函数只含 $z$（或含 $z$ 的因子可分离）且截面面积易求时，内层二重积分 $=$ 截面面积 $\times$ 该因子。

\corepoint{两法的差别只在"先算哪个"：投影法是\textbf{先一后二}，要求沿 $z$ 方向穿入穿出各一次；截面法是\textbf{先二后一}，要求截面形状随 $z$ 变化简单。判据：被积函数中 $z$ 的因子能否提出内层——能提就优先截面法。}

\subsection{柱面坐标}

设点 $M$ 在 $xOy$ 面上投影 $P$ 的极坐标为 $(\rho,\theta)$，则 $(\rho,\theta,z)$ 为 $M$ 的\textbf{柱面坐标}，$0\le\rho<+\infty$，$0\le\theta\le2\pi$，$-\infty<z<+\infty$。三组坐标面：$\rho=$ 常数（以 $z$ 轴为轴的圆柱面）、$\theta=$ 常数（过 $z$ 轴的半平面）、$z=$ 常数（平行于 $xOy$ 面的平面）。关系式与体积元素：
\keyformula{x=\rho\cos\theta,\quad y=\rho\sin\theta,\quad z=z;\qquad \mathrm{d}V=\rho\,\mathrm{d}\rho\,\mathrm{d}\theta\,\mathrm{d}z}
\keyformula{\iiint_\Omega f\,\mathrm{d}V=\iiint_{\Omega'}f(\rho\cos\theta,\rho\sin\theta,z)\,\rho\,\mathrm{d}\rho\,\mathrm{d}\theta\,\mathrm{d}z}
再化为三次积分，积分限按 $\rho,\theta,z$ 在 $\Omega$ 中的变化范围确定。

\corepoint{柱面坐标适合"绕 $z$ 轴旋转"的区域（圆柱、圆锥、旋转抛物面、圆环）与被积函数含 $x^2+y^2$ 的情形。$\rho$ 的限与二重积分的极坐标一致：先看投影区域张多大角，再看沿射线 $\rho$ 从哪进哪出；别忘了因子 $\rho$。}

\subsection{球面坐标}

设 $r$ 为原点 $O$ 与 $M$ 的距离，$\varphi$ 为有向线段 $\overrightarrow{OM}$ 与 $z$ 轴正向的夹角，$\theta$ 为从 $z$ 轴正向来看自 $x$ 轴按逆时针转到 $\overrightarrow{OP}$ 的角（$P$ 为 $M$ 在 $xOy$ 面上的投影），则 $(r,\varphi,\theta)$ 为 $M$ 的\textbf{球面坐标}，$0\le r<+\infty$，$0\le\varphi\le\pi$，$0\le\theta\le2\pi$。三组坐标面：$r=$ 常数（以原点为球心的球面）、$\varphi=$ 常数（以原点为顶点、$z$ 轴为轴的圆锥面）、$\theta=$ 常数（过 $z$ 轴的半平面）。
\keyformula{x=r\sin\varphi\cos\theta,\quad y=r\sin\varphi\sin\theta,\quad z=r\cos\varphi}
\keyformula{\mathrm{d}V=r^{2}\sin\varphi\,\mathrm{d}r\,\mathrm{d}\varphi\,\mathrm{d}\theta}
\keyformula{\iiint_\Omega f\,\mathrm{d}V=\iiint_{\Omega'}f(r\sin\varphi\cos\theta,\ r\sin\varphi\sin\theta,\ r\cos\varphi)\,r^{2}\sin\varphi\,\mathrm{d}r\,\mathrm{d}\varphi\,\mathrm{d}\theta}
一般化为先对 $r$、再对 $\varphi$、最后对 $\theta$ 的三次积分：$0\le\theta\le2\pi$，$0\le\varphi\le\pi$（只在上半球时 $0\le\varphi\le\frac{\pi}{2}$），$r$ 的限由"自原点出发的射线穿过 $\Omega$"确定。

\corepoint{两个最易错处：(1) 体积元素是 $r^2\sin\varphi$，不是 $\cos\varphi$、也不是 $\sin\theta$；(2) $\varphi$ 是\textbf{与 $z$ 轴正向的夹角}，不是与 $xOy$ 面的夹角——上半球对应 $\varphi\in[0,\frac{\pi}{2}]$，圆锥面 $z=\sqrt{x^2+y^2}$ 对应 $\varphi=\frac{\pi}{4}$。记不住就用 $z=r\cos\varphi$ 现推。}

\subsection{三种坐标的选择}

\begin{center}
\begin{tabular}{@{}lll@{}}
\toprule
坐标 & 适用的积分区域 & 体积元素 \\
\midrule
直角坐标 & 长方体、四面体、由平面与柱面围成的可分层区域 & $\mathrm{d}x\,\mathrm{d}y\,\mathrm{d}z$ \\
柱面坐标 & 圆柱、圆锥、旋转抛物面、绕 $z$ 轴的旋转体 & $\rho\,\mathrm{d}\rho\,\mathrm{d}\theta\,\mathrm{d}z$ \\
球面坐标 & 球体、球壳、球与锥面所围的"冰淇淋"形 & $r^{2}\sin\varphi\,\mathrm{d}r\,\mathrm{d}\varphi\,\mathrm{d}\theta$ \\
\bottomrule
\end{tabular}
\end{center}

选择原则：先看\textbf{边界曲面}（球用球面坐标、柱面用柱面坐标），再看\textbf{被积函数}含 $x^2+y^2$（柱面坐标）还是 $x^2+y^2+z^2$（球面坐标）；冲突时以"积分限好定"为准。

\section{重积分的应用}
\subsection{元素法}

若所求量 $U$ (1) 对积分区域具有可加性，(2) 在区域内某点的很小邻域上其部分量可用 $f(M)\,\mathrm{d}\Omega$ 近似（$\mathrm{d}\Omega$ 为面积或体积元素），则 $U=\iint_Df\,\mathrm{d}\sigma$ 或 $U=\iiint_\Omega f\,\mathrm{d}V$。关键是写出\textbf{正确的微元}：几何量用面积（体积）元素，物理量用"密度 $\times$ 元素"。

\subsection{曲面的面积}

设曲面 $S:z=f(x,y)$，$f$ 在 $D$（$S$ 在 $xOy$ 面上的投影区域）有连续偏导。$S$ 上小片面积 $\mathrm{d}A$ 与其投影 $\mathrm{d}\sigma$ 之比为 $\frac{1}{\cos\gamma}$（$\gamma$ 为 $M$ 处法线与 $z$ 轴夹角），而 $\cos\gamma=\frac{1}{\sqrt{1+f_x^2+f_y^2}}$，故\textbf{曲面面积元素} $\mathrm{d}A=\sqrt{1+f_x^2+f_y^2}\,\mathrm{d}\sigma$，
\keyformula{A=\iint_D\sqrt{1+\left(\frac{\partial z}{\partial x}\right)^{2}+\left(\frac{\partial z}{\partial y}\right)^{2}}\ \mathrm{d}x\,\mathrm{d}y}
若 $S$：$x=g(y,z)$ 或 $y=h(z,x)$，可分别投影到 $yOz$ 面、$zOx$ 面，把根号内换成相应两偏导的平方和。

\corepoint{公式要求被积函数在 $D$ 上连续。若 $S$ 在投影区域边界处法线趋于与 $z$ 轴垂直（$\cos\gamma\to0$，如球面的"极点"、柱面的母线），被积函数在 $D$ 上无界，不能直接套用，应先取内接区域算出一部分再取极限。}

\subsection{质心}

平面薄片（面密度 $\mu$ 连续）的质量与质心：
\keyformula{M=\iint_D\mu\,\mathrm{d}\sigma,\qquad \bar x=\frac{1}{M}\iint_Dx\mu\,\mathrm{d}\sigma,\qquad \bar y=\frac{1}{M}\iint_Dy\mu\,\mathrm{d}\sigma}
空间物体（体密度 $\mu$ 连续）：$M=\iiint_\Omega\mu\,\mathrm{d}V$，$\bar x=\frac{1}{M}\iiint_\Omega x\mu\,\mathrm{d}V$，$\bar y,\bar z$ 同理。$\mu$ 为常数时质心即\textbf{形心}，$\mu$ 可约去。

\subsection{转动惯量}

平面薄片对 $x$ 轴、$y$ 轴、原点的转动惯量：
\keyformula{I_x=\iint_Dy^{2}\mu\,\mathrm{d}\sigma,\qquad I_y=\iint_Dx^{2}\mu\,\mathrm{d}\sigma,\qquad I_O=I_x+I_y=\iint_D(x^{2}+y^{2})\mu\,\mathrm{d}\sigma}
空间物体：
\[
I_x=\iiint_\Omega(y^2+z^2)\mu\,\mathrm{d}V,\quad
I_y=\iiint_\Omega(z^2+x^2)\mu\,\mathrm{d}V,\quad
I_z=\iiint_\Omega(x^2+y^2)\mu\,\mathrm{d}V,\quad
I_O=I_x+I_y+I_z .
\]
记忆要点：\textbf{对哪根轴取转动惯量，微元里的"距离平方"就是另外两个坐标的平方和}。

\subsection{引力}

设空间物体占有 $\Omega$，体密度 $\mu$ 连续，则它对位于 $M_0(x_0,y_0,z_0)$、质量为 $m$ 的质点的引力 $\boldsymbol F=(F_x,F_y,F_z)$ 为
\[
F_x=km\iiint_\Omega\frac{\mu(x,y,z)(x-x_0)}{r^{3}}\,\mathrm{d}V ,
\]
$F_y,F_z$ 只需把分子换成 $y-y_0$、$z-z_0$，其中 $r=\sqrt{(x-x_0)^2+(y-y_0)^2+(z-z_0)^2}$，$k$ 为引力常数。

\corepoint{分母是 $r^{\mathbf{3}}$ 而非 $r^2$：微元对单位质量的引力大小为 $\frac{k\mu\,\mathrm{d}V}{r^{2}}$，再乘方向余弦 $\frac{x-x_0}{r}$，合起来即 $\frac{x-x_0}{r^3}$。平面薄片对质点的引力把三重积分换成二重积分、$\mathrm{d}V$ 换成 $\mathrm{d}\sigma$，$r$ 仍按空间距离算。}

\section{含参变量的积分}

设 $f(x,y)$ 在矩形 $R=[a,b]\times[c,d]$ 上连续，则 $\Phi(y)=\int_a^bf(x,y)\,\mathrm{d}x$ 确定了 $y$ 在 $[c,d]$ 上的一个函数，称为\textbf{含参变量的积分}（$y$ 为参变量）。

\textbf{定理 1（连续性）}\quad 若 $f(x,y)$ 在 $R$ 上连续，则 $\Phi(y)$ 在 $[c,d]$ 上连续，即 $\lim\limits_{y\to y_0}\int_a^bf(x,y)\,\mathrm{d}x=\int_a^bf(x,y_0)\,\mathrm{d}x$（极限与积分可交换）。

\textbf{定理 2（积分次序交换）}\quad 若 $f$ 在 $R$ 上连续，则
\[
\int_c^d\mathrm{d}y\int_a^bf(x,y)\,\mathrm{d}x=\int_a^b\mathrm{d}x\int_c^df(x,y)\,\mathrm{d}y
\]
（即矩形区域上两种次序的二次积分相等）。

\textbf{定理 3（可导性·莱布尼茨公式）}\quad 若 $f$ 及其偏导数 $f_y$ 都在 $R$ 上连续，则 $\Phi(y)$ 在 $[c,d]$ 上可微，且
\keyformula{\Phi'(y)=\frac{\mathrm{d}}{\mathrm{d}y}\int_a^bf(x,y)\,\mathrm{d}x=\int_a^b\frac{\partial f(x,y)}{\partial y}\,\mathrm{d}x}
即\textbf{积分与求导可交换次序}。

\textbf{定理 4（积分限也含参变量）}\quad 若 $f$ 在 $R$ 上连续，$\alpha(y),\beta(y)$ 可微且 $a\le\alpha,\beta\le b$，则 $\Phi(y)=\int_{\alpha(y)}^{\beta(y)}f(x,y)\,\mathrm{d}x$ 可微，且
\[
\Phi'(y)=\int_{\alpha(y)}^{\beta(y)}f_y(x,y)\,\mathrm{d}x+f(\beta(y),y)\beta'(y)-f(\alpha(y),y)\alpha'(y).
\]

\textbf{常用方法}：(1) \textbf{对参变量微分法}——把待求定积分看作某含参积分的特殊值，用莱布尼茨公式求 $\Phi'(y)$，再积分回 $\Phi(y)$ 并代特殊值；(2) \textbf{先积分后求导}——先把 $\Phi(y)$ 算出再求导；(3) \textbf{交换积分次序}——把一元定积分凑成二重积分再换序。

\corepoint{用莱布尼茨公式前必须核对条件：$f$ 与 $f_y$ 都要在\textbf{闭矩形}上连续。积分限含参变量时别漏后两项。用"对参数微分法"求定积分时，引入参数要使 $\Phi(y_0)$ 好算，且全过程参数不得越出连续性所在的矩形。}
 引入，不含 document 环境

\section{二重积分的概念与性质}
\subsection{引例与定义}

\textbf{引例一：曲顶柱体的体积。}底为 $xOy$ 面上的闭区域 $D$，顶为曲面 $z=f(x,y)$（$f\ge0$ 连续），高随点而变，"体积 $=$ 高 $\times$ 底面积"失效。用曲线网把 $D$ 分成 $n$ 个小闭区域 $\Delta\sigma_i$（既表示小块，也表示其面积），细曲顶柱体近似为以 $f(\xi_i,\eta_i)$ 为高的平顶柱体，求和后令最大直径 $\lambda\to0$ 取极限：$V=\lim\limits_{\lambda\to0}\sum\limits_{i=1}^{n}f(\xi_i,\eta_i)\Delta\sigma_i$。

\textbf{引例二：平面薄片的质量。}薄片占有 $D$，面密度 $\mu(x,y)>0$ 连续，同样"分割—近似—求和—取极限"，得 $m=\lim\limits_{\lambda\to0}\sum\limits_{i=1}^{n}\mu(\xi_i,\eta_i)\Delta\sigma_i$。

两问题的实际意义不同，但所求量都归结为同一形式的和的极限，抽去具体意义即得

\textbf{定义（二重积分）}\quad 设 $f(x,y)$ 是有界闭区域 $D$ 上的有界函数。将 $D$ 任意分成 $n$ 个小闭区域 $\Delta\sigma_i$，在每个 $\Delta\sigma_i$ 上任取一点 $(\xi_i,\eta_i)$，作和 $\sum\limits_{i=1}^{n}f(\xi_i,\eta_i)\Delta\sigma_i$。若当 $\lambda\to0$ 时该和的极限总存在，且与 $D$ 的分法及点 $(\xi_i,\eta_i)$ 的取法无关，则称此极限为 $f(x,y)$ 在 $D$ 上的\textbf{二重积分}，记作
\[
\iint_D f(x,y)\,\mathrm{d}\sigma=\lim_{\lambda\to0}\sum_{i=1}^{n}f(\xi_i,\eta_i)\,\Delta\sigma_i ,
\]
其中 $f(x,y)$ 为\textbf{被积函数}，$f(x,y)\mathrm{d}\sigma$ 为\textbf{被积表达式}，$\mathrm{d}\sigma$ 为\textbf{面积元素}，$x,y$ 为\textbf{积分变量}，$D$ 为\textbf{积分区域}。

\begin{itemize}
  \item 直角坐标中用平行于坐标轴的直线网划分 $D$，除含边界点的小块外都是矩形，$\Delta\sigma=\Delta x\,\Delta y$，故 $\mathrm{d}\sigma=\mathrm{d}x\,\mathrm{d}y$，二重积分也记作 $\iint_Df(x,y)\,\mathrm{d}x\,\mathrm{d}y$。
  \item \textbf{可积性}：$f(x,y)$ 在 $D$ 上\textbf{连续}时积分和极限必存在，即必可积（可积的必要条件是被积函数有界）。教材以后总假定可积。
  \item \textbf{几何意义}：$f\ge0$ 时积分值是以 $D$ 为底、$z=f(x,y)$ 为顶的曲顶柱体体积；$f$ 变号时，积分值等于 $xOy$ 面上方柱体体积减去下方柱体体积之差。
\end{itemize}

\subsection{二重积分的性质}

设所涉积分都存在，$\sigma$ 为 $D$ 的面积。

\begin{enumerate}
  \item \textbf{线性性质}\quad $\iint_D[\alpha f+\beta g]\,\mathrm{d}\sigma=\alpha\iint_Df\,\mathrm{d}\sigma+\beta\iint_Dg\,\mathrm{d}\sigma$（$\alpha,\beta$ 为常数）。
  \item \textbf{对积分区域的可加性}\quad 若 $D$ 被有限条曲线分为有限个部分闭区域（如 $D=D_1\cup D_2$，无公共内点），则 $\iint_Df\,\mathrm{d}\sigma=\iint_{D_1}f\,\mathrm{d}\sigma+\iint_{D_2}f\,\mathrm{d}\sigma$。
  \item \textbf{面积性质}\quad 若在 $D$ 上 $f\equiv1$，则 $\iint_D1\,\mathrm{d}\sigma=\iint_D\mathrm{d}\sigma=\sigma$（高为 $1$ 的平顶柱体体积在数值上等于底面积）。
  \item \textbf{比较性质}\quad 若在 $D$ 上 $f\le g$，则 $\iint_Df\,\mathrm{d}\sigma\le\iint_Dg\,\mathrm{d}\sigma$；特别地 $\left|\iint_Df\,\mathrm{d}\sigma\right|\le\iint_D|f|\,\mathrm{d}\sigma$。
  \item \textbf{估值性质}\quad 设 $M,m$ 为 $f$ 在 $D$ 上的最大值、最小值，则 $m\sigma\le\iint_Df\,\mathrm{d}\sigma\le M\sigma$。
  \item \textbf{中值定理}\quad 设 $f$ 在 $D$ 上连续，则 $D$ 上至少存在一点 $(\xi,\eta)$ 使 $\iint_Df(x,y)\,\mathrm{d}\sigma=f(\xi,\eta)\sigma$。（证：由性质 5 得 $m\le\frac{1}{\sigma}\iint_Df\,\mathrm{d}\sigma\le M$，再用连续函数的介值定理。）
\end{enumerate}

\corepoint{性质 5、6 中的"$\sigma$ 是 $D$ 的面积"不能漏：直角坐标下 $\sigma=\iint_D\mathrm{d}x\mathrm{d}y$，极坐标下 $\sigma=\iint_D\rho\,\mathrm{d}\rho\mathrm{d}\theta$，换元后面积要重算，不能把 $\Delta x\Delta y$ 当面积。另外"$f\le g\Rightarrow$ 积分值 $\le$"只对\textbf{同一积分区域}成立。}

\subsection{对称性简化（常用）}

\begin{itemize}
  \item $D$ 关于 $y$ 轴对称时：
  \[
  \iint_Df(x,y)\,\mathrm{d}\sigma=
  \begin{cases}
  0, & f(-x,y)=-f(x,y),\\[2pt]
  2\displaystyle\iint_{D_1}f(x,y)\,\mathrm{d}\sigma, & f(-x,y)=f(x,y),
  \end{cases}
  \]
  $D_1$ 为 $D$ 中 $x\ge0$ 的部分。
  \item $D$ 关于 $x$ 轴对称时把 $x,y$ 互换；若 $D$ 关于两轴均对称且 $f$ 关于 $x,y$ 都是偶函数，则 $\iint_Df\,\mathrm{d}\sigma=4\iint_{D_1}f\,\mathrm{d}\sigma$（$D_1$ 为第一象限部分）；$D$ 关于原点对称且 $f(-x,-y)=f(x,y)$ 时 $\iint_Df\,\mathrm{d}\sigma=2\iint_{D_{y\ge0}}f\,\mathrm{d}\sigma$。
\end{itemize}

\corepoint{用对称性有\textbf{两个前提}：(1) 积分\textbf{区域}对称；(2) 被积函数在该对称下有确定的奇偶性，缺一不可。判断奇偶性时只有一个变量变号，其余当常数。若 $D$ 关于 $y$ 轴对称而 $f$ 关于 $x$ 为奇函数，则积分为 $0$，这一句常能省掉整段计算。}

\section{二重积分的计算法}

按定义计算只对少数特别简单的被积函数与区域可行。基本方法是把二重积分化为\textbf{两次定积分}（累次积分）。

\subsection{利用直角坐标计算}

取 $\mathrm{d}\sigma=\mathrm{d}x\,\mathrm{d}y$。

\textbf{$X$ 型区域（先对 $y$、后对 $x$）}\quad 若 $D$ 可用不等式 $a\le x\le b$，$\varphi_1(x)\le y\le\varphi_2(x)$ 表示（$\varphi_1\le\varphi_2$ 在 $[a,b]$ 上连续），则
\keyformula{\iint_Df(x,y)\,\mathrm{d}x\,\mathrm{d}y=\int_a^b\mathrm{d}x\int_{\varphi_1(x)}^{\varphi_2(x)}f(x,y)\,\mathrm{d}y}
特征：穿过 $D$ 内部且\textbf{平行于 $y$ 轴}的直线与 $D$ 的边界相交不多于两点。

\textbf{$Y$ 型区域（先对 $x$、后对 $y$）}\quad 若 $D$ 可用 $c\le y\le d$，$\psi_1(y)\le x\le\psi_2(y)$ 表示，则
\keyformula{\iint_Df(x,y)\,\mathrm{d}x\,\mathrm{d}y=\int_c^d\mathrm{d}y\int_{\psi_1(y)}^{\psi_2(y)}f(x,y)\,\mathrm{d}x}
特征：穿过 $D$ 内部且\textbf{平行于 $x$ 轴}的直线与边界相交不多于两点。

两种情形要处理：

\begin{itemize}
  \item 若既有平行于 $y$ 轴的直线与边界相交多于两点、又有平行于 $x$ 轴的直线相交多于两点，则 $D$ 既非 $X$ 型也非 $Y$ 型，须把 $D$ 分成几块，每块为 $X$ 型或 $Y$ 型，再由性质 2 相加。
  \item 若 $D$ 既是 $X$ 型又是 $Y$ 型，则两种次序的二次积分相等，都等于同一个二重积分（这正是累次积分可交换次序的依据）。
\end{itemize}

\textbf{化累次积分的步骤}

\begin{enumerate}
  \item \textbf{画区域}：画出 $D$，求出边界曲线的全部交点。
  \item \textbf{定类型}：判断 $X$ 型、$Y$ 型还是需分块。
  \item \textbf{定限}：外层积分限必为\textbf{常数}，内层积分限是外层变量的函数。
  \item \textbf{算内层}：把外层变量当常数，得到外层变量的函数。
  \item \textbf{算外层}：对剩下的变量作定积分。
\end{enumerate}

\corepoint{定限是唯一的关键，口诀"先画图、后定限，外层常数、内层函数"。常见错误：外层限写成函数；上下限写反；分块时漏块；先对 $y$ 积时把含 $x$ 的因子留在内层（应作为常数提出）。}

\corepoint{\textbf{选择次序}要同时看区域形状与被积函数：既要 $D$ 少分块，又要内层能积出来。含 $\mathrm{e}^{y^2}$、$\frac{\sin y}{y}$、$\sqrt{1+y^3}$ 之类的积分对 $y$ 积不出来，必须\textbf{交换次序}：先由原累次积分还原 $D$ 的不等式组，画图，再按新次序重新定限。这常常是唯一出路。}

\subsection{利用极坐标计算}

当 $D$ 的边界用极坐标表示方便（圆、圆环、扇形）且被积函数含 $x^2+y^2$、$\frac{y}{x}$ 时用极坐标。取 $x=\rho\cos\theta$，$y=\rho\sin\theta$，两族曲线围成的小块面积 $\approx\rho\,\mathrm{d}\rho\,\mathrm{d}\theta$，故\textbf{极坐标下的面积元素} $\mathrm{d}\sigma=\rho\,\mathrm{d}\rho\,\mathrm{d}\theta$，
\keyformula{\iint_Df(x,y)\,\mathrm{d}x\,\mathrm{d}y=\iint_Df(\rho\cos\theta,\rho\sin\theta)\,\rho\,\mathrm{d}\rho\,\mathrm{d}\theta}
即：$x,y$ 分别换成 $\rho\cos\theta,\rho\sin\theta$，$\mathrm{d}x\mathrm{d}y$ 换成 $\rho\,\mathrm{d}\rho\mathrm{d}\theta$（\textbf{这个 $\rho$ 不能漏}）。

设 $D$ 为 $\rho_1(\theta)\le\rho\le\rho_2(\theta)$，$\alpha\le\theta\le\beta$，则
\keyformula{\iint_Df(\rho\cos\theta,\rho\sin\theta)\,\rho\,\mathrm{d}\rho\,\mathrm{d}\theta=\int_\alpha^\beta\mathrm{d}\theta\int_{\rho_1(\theta)}^{\rho_2(\theta)}f(\rho\cos\theta,\rho\sin\theta)\,\rho\,\mathrm{d}\rho}
两种特例：极点在 $D$ 的\textbf{边界}上（曲边扇形）时 $\rho_1\equiv0$，$\rho$ 的上限为 $\rho(\theta)$；极点在 $D$ 的\textbf{内部}时 $0\le\rho\le\rho(\theta)$，$0\le\theta\le2\pi$。由性质 3 得曲边扇形的面积 $\sigma=\frac{1}{2}\int_\alpha^\beta[\rho_2^2(\theta)-\rho_1^2(\theta)]\,\mathrm{d}\theta$。

\corepoint{极坐标三个高频错误：(1) 漏掉面积元素里的 $\rho$；(2) 不管区域张多大角，一律写 $[0,2\pi]$；(3) 内外层限的位置搞反。另需记准两个常用转换：$x^2+y^2=2ax$ 即 $\rho=2a\cos\theta$（$\theta\in[-\frac{\pi}{2},\frac{\pi}{2}]$），$x^2+y^2=2ay$ 即 $\rho=2a\sin\theta$（$\theta\in[0,\pi]$）。}

\subsection{二重积分的换元法}

设 $f(x,y)$ 在 $D$ 上连续，变换 $x=x(u,v)$，$y=y(u,v)$ 把 $uOv$ 面上的闭区域 $D'$ 一一对应地映成 $D$，$x,y$ 在 $D'$ 上有一阶连续偏导数，其\textbf{雅可比行列式}
\[
J(u,v)=\frac{\partial(x,y)}{\partial(u,v)}
=\begin{vmatrix}
\dfrac{\partial x}{\partial u} & \dfrac{\partial x}{\partial v}\\[8pt]
\dfrac{\partial y}{\partial u} & \dfrac{\partial y}{\partial v}
\end{vmatrix}\ne0 ,
\]
则
\keyformula{\iint_Df(x,y)\,\mathrm{d}x\,\mathrm{d}y=\iint_{D'}f\bigl(x(u,v),y(u,v)\bigr)\,\bigl|J(u,v)\bigr|\,\mathrm{d}u\,\mathrm{d}v}
极坐标 $x=\rho\cos\theta,y=\rho\sin\theta$ 是其特例，$|J|=\rho$。常用\textbf{广义极坐标}（椭圆域 $\frac{x^2}{a^2}+\frac{y^2}{b^2}\le1$）：$x=a\rho\cos\theta$，$y=b\rho\sin\theta$，$|J|=ab\rho$。

\corepoint{公式里是 $|J|$ 的\textbf{绝对值}；换元后区域一并换成 $D'$，积分限须按 $D'$ 重定；变换须一一对应（$J\ne0$ 只保证局部一一），否则分块。}

\section{三重积分}
\subsection{概念}

\textbf{定义}\quad 设 $f(x,y,z)$ 是空间有界闭区域 $\Omega$ 上的有界函数。把 $\Omega$ 任意分成 $n$ 个小闭区域 $\Delta V_i$（也表示其体积），任取 $(\xi_i,\eta_i,\zeta_i)$ 作和 $\sum f(\xi_i,\eta_i,\zeta_i)\Delta V_i$。若 $\lambda\to0$（各小块直径最大值）时极限总存在且与分法、取法无关，则称之为 $f$ 在 $\Omega$ 上的\textbf{三重积分}：
\[
\iiint_\Omega f(x,y,z)\,\mathrm{d}V=\lim_{\lambda\to0}\sum_{i=1}^{n}f(\xi_i,\eta_i,\zeta_i)\,\Delta V_i ,
\]
$\mathrm{d}V$ 为\textbf{体积元素}；直角坐标下 $\mathrm{d}V=\mathrm{d}x\,\mathrm{d}y\,\mathrm{d}z$。$f$ 连续时必可积；$f\ge0$ 时可解释为空间物体的质量（$f$ 为体密度）。

性质与二重积分完全平行：线性、对区域可加、$\iiint_\Omega\mathrm{d}V=V$、比较、估值、中值定理 $\iiint_\Omega f\,\mathrm{d}V=f(\xi,\eta,\zeta)V$。

\subsection{直角坐标：投影法与截面法}

\textbf{投影法（先一后二）}\quad 设平行于 $z$ 轴且穿过 $\Omega$ 内部的直线与边界曲面 $S$ 相交不多于两点。把 $\Omega$ 投影到 $xOy$ 面得 $D_{xy}$，$S$ 被分成下、上两部分 $S_1:z=z_1(x,y)$，$S_2:z=z_2(x,y)$（$z_1\le z_2$ 均在 $D_{xy}$ 上连续），则
\keyformula{\iiint_\Omega f\,\mathrm{d}V=\iint_{D_{xy}}\mathrm{d}x\,\mathrm{d}y\int_{z_1(x,y)}^{z_2(x,y)}f(x,y,z)\,\mathrm{d}z}
即先对 $z$、再对 $x,y$ 的三次积分。若平行于 $x$ 轴或 $y$ 轴的直线与 $S$ 相交不多于两点，也可投影到 $yOz$ 面或 $zOx$ 面，化为其他顺序的三次积分；交点多于两点时把 $\Omega$ 分块相加。

\textbf{截面法（先二后一）}\quad 设 $\Omega=\{(x,y,z)\mid c\le z\le d,\ (x,y)\in D_z\}$，其中 $D_z$ 是用竖坐标 $z$ 的平面截 $\Omega$ 所得的平面闭区域，则
\keyformula{\iiint_\Omega f\,\mathrm{d}V=\int_c^d\mathrm{d}z\iint_{D_z}f(x,y,z)\,\mathrm{d}x\,\mathrm{d}y}
当被积函数只含 $z$（或含 $z$ 的因子可分离）且截面面积易求时，内层二重积分 $=$ 截面面积 $\times$ 该因子。

\corepoint{两法的差别只在"先算哪个"：投影法是\textbf{先一后二}，要求沿 $z$ 方向穿入穿出各一次；截面法是\textbf{先二后一}，要求截面形状随 $z$ 变化简单。判据：被积函数中 $z$ 的因子能否提出内层——能提就优先截面法。}

\subsection{柱面坐标}

设点 $M$ 在 $xOy$ 面上投影 $P$ 的极坐标为 $(\rho,\theta)$，则 $(\rho,\theta,z)$ 为 $M$ 的\textbf{柱面坐标}，$0\le\rho<+\infty$，$0\le\theta\le2\pi$，$-\infty<z<+\infty$。三组坐标面：$\rho=$ 常数（以 $z$ 轴为轴的圆柱面）、$\theta=$ 常数（过 $z$ 轴的半平面）、$z=$ 常数（平行于 $xOy$ 面的平面）。关系式与体积元素：
\keyformula{x=\rho\cos\theta,\quad y=\rho\sin\theta,\quad z=z;\qquad \mathrm{d}V=\rho\,\mathrm{d}\rho\,\mathrm{d}\theta\,\mathrm{d}z}
\keyformula{\iiint_\Omega f\,\mathrm{d}V=\iiint_{\Omega'}f(\rho\cos\theta,\rho\sin\theta,z)\,\rho\,\mathrm{d}\rho\,\mathrm{d}\theta\,\mathrm{d}z}
再化为三次积分，积分限按 $\rho,\theta,z$ 在 $\Omega$ 中的变化范围确定。

\corepoint{柱面坐标适合"绕 $z$ 轴旋转"的区域（圆柱、圆锥、旋转抛物面、圆环）与被积函数含 $x^2+y^2$ 的情形。$\rho$ 的限与二重积分的极坐标一致：先看投影区域张多大角，再看沿射线 $\rho$ 从哪进哪出；别忘了因子 $\rho$。}

\subsection{球面坐标}

设 $r$ 为原点 $O$ 与 $M$ 的距离，$\varphi$ 为有向线段 $\overrightarrow{OM}$ 与 $z$ 轴正向的夹角，$\theta$ 为从 $z$ 轴正向来看自 $x$ 轴按逆时针转到 $\overrightarrow{OP}$ 的角（$P$ 为 $M$ 在 $xOy$ 面上的投影），则 $(r,\varphi,\theta)$ 为 $M$ 的\textbf{球面坐标}，$0\le r<+\infty$，$0\le\varphi\le\pi$，$0\le\theta\le2\pi$。三组坐标面：$r=$ 常数（以原点为球心的球面）、$\varphi=$ 常数（以原点为顶点、$z$ 轴为轴的圆锥面）、$\theta=$ 常数（过 $z$ 轴的半平面）。
\keyformula{x=r\sin\varphi\cos\theta,\quad y=r\sin\varphi\sin\theta,\quad z=r\cos\varphi}
\keyformula{\mathrm{d}V=r^{2}\sin\varphi\,\mathrm{d}r\,\mathrm{d}\varphi\,\mathrm{d}\theta}
\keyformula{\iiint_\Omega f\,\mathrm{d}V=\iiint_{\Omega'}f(r\sin\varphi\cos\theta,\ r\sin\varphi\sin\theta,\ r\cos\varphi)\,r^{2}\sin\varphi\,\mathrm{d}r\,\mathrm{d}\varphi\,\mathrm{d}\theta}
一般化为先对 $r$、再对 $\varphi$、最后对 $\theta$ 的三次积分：$0\le\theta\le2\pi$，$0\le\varphi\le\pi$（只在上半球时 $0\le\varphi\le\frac{\pi}{2}$），$r$ 的限由"自原点出发的射线穿过 $\Omega$"确定。

\corepoint{两个最易错处：(1) 体积元素是 $r^2\sin\varphi$，不是 $\cos\varphi$、也不是 $\sin\theta$；(2) $\varphi$ 是\textbf{与 $z$ 轴正向的夹角}，不是与 $xOy$ 面的夹角——上半球对应 $\varphi\in[0,\frac{\pi}{2}]$，圆锥面 $z=\sqrt{x^2+y^2}$ 对应 $\varphi=\frac{\pi}{4}$。记不住就用 $z=r\cos\varphi$ 现推。}

\subsection{三种坐标的选择}

\begin{center}
\begin{tabular}{@{}lll@{}}
\toprule
坐标 & 适用的积分区域 & 体积元素 \\
\midrule
直角坐标 & 长方体、四面体、由平面与柱面围成的可分层区域 & $\mathrm{d}x\,\mathrm{d}y\,\mathrm{d}z$ \\
柱面坐标 & 圆柱、圆锥、旋转抛物面、绕 $z$ 轴的旋转体 & $\rho\,\mathrm{d}\rho\,\mathrm{d}\theta\,\mathrm{d}z$ \\
球面坐标 & 球体、球壳、球与锥面所围的"冰淇淋"形 & $r^{2}\sin\varphi\,\mathrm{d}r\,\mathrm{d}\varphi\,\mathrm{d}\theta$ \\
\bottomrule
\end{tabular}
\end{center}

选择原则：先看\textbf{边界曲面}（球用球面坐标、柱面用柱面坐标），再看\textbf{被积函数}含 $x^2+y^2$（柱面坐标）还是 $x^2+y^2+z^2$（球面坐标）；冲突时以"积分限好定"为准。

\section{重积分的应用}
\subsection{元素法}

若所求量 $U$ (1) 对积分区域具有可加性，(2) 在区域内某点的很小邻域上其部分量可用 $f(M)\,\mathrm{d}\Omega$ 近似（$\mathrm{d}\Omega$ 为面积或体积元素），则 $U=\iint_Df\,\mathrm{d}\sigma$ 或 $U=\iiint_\Omega f\,\mathrm{d}V$。关键是写出\textbf{正确的微元}：几何量用面积（体积）元素，物理量用"密度 $\times$ 元素"。

\subsection{曲面的面积}

设曲面 $S:z=f(x,y)$，$f$ 在 $D$（$S$ 在 $xOy$ 面上的投影区域）有连续偏导。$S$ 上小片面积 $\mathrm{d}A$ 与其投影 $\mathrm{d}\sigma$ 之比为 $\frac{1}{\cos\gamma}$（$\gamma$ 为 $M$ 处法线与 $z$ 轴夹角），而 $\cos\gamma=\frac{1}{\sqrt{1+f_x^2+f_y^2}}$，故\textbf{曲面面积元素} $\mathrm{d}A=\sqrt{1+f_x^2+f_y^2}\,\mathrm{d}\sigma$，
\keyformula{A=\iint_D\sqrt{1+\left(\frac{\partial z}{\partial x}\right)^{2}+\left(\frac{\partial z}{\partial y}\right)^{2}}\ \mathrm{d}x\,\mathrm{d}y}
若 $S$：$x=g(y,z)$ 或 $y=h(z,x)$，可分别投影到 $yOz$ 面、$zOx$ 面，把根号内换成相应两偏导的平方和。

\corepoint{公式要求被积函数在 $D$ 上连续。若 $S$ 在投影区域边界处法线趋于与 $z$ 轴垂直（$\cos\gamma\to0$，如球面的"极点"、柱面的母线），被积函数在 $D$ 上无界，不能直接套用，应先取内接区域算出一部分再取极限。}

\subsection{质心}

平面薄片（面密度 $\mu$ 连续）的质量与质心：
\keyformula{M=\iint_D\mu\,\mathrm{d}\sigma,\qquad \bar x=\frac{1}{M}\iint_Dx\mu\,\mathrm{d}\sigma,\qquad \bar y=\frac{1}{M}\iint_Dy\mu\,\mathrm{d}\sigma}
空间物体（体密度 $\mu$ 连续）：$M=\iiint_\Omega\mu\,\mathrm{d}V$，$\bar x=\frac{1}{M}\iiint_\Omega x\mu\,\mathrm{d}V$，$\bar y,\bar z$ 同理。$\mu$ 为常数时质心即\textbf{形心}，$\mu$ 可约去。

\subsection{转动惯量}

平面薄片对 $x$ 轴、$y$ 轴、原点的转动惯量：
\keyformula{I_x=\iint_Dy^{2}\mu\,\mathrm{d}\sigma,\qquad I_y=\iint_Dx^{2}\mu\,\mathrm{d}\sigma,\qquad I_O=I_x+I_y=\iint_D(x^{2}+y^{2})\mu\,\mathrm{d}\sigma}
空间物体：
\[
I_x=\iiint_\Omega(y^2+z^2)\mu\,\mathrm{d}V,\quad
I_y=\iiint_\Omega(z^2+x^2)\mu\,\mathrm{d}V,\quad
I_z=\iiint_\Omega(x^2+y^2)\mu\,\mathrm{d}V,\quad
I_O=I_x+I_y+I_z .
\]
记忆要点：\textbf{对哪根轴取转动惯量，微元里的"距离平方"就是另外两个坐标的平方和}。

\subsection{引力}

设空间物体占有 $\Omega$，体密度 $\mu$ 连续，则它对位于 $M_0(x_0,y_0,z_0)$、质量为 $m$ 的质点的引力 $\boldsymbol F=(F_x,F_y,F_z)$ 为
\[
F_x=km\iiint_\Omega\frac{\mu(x,y,z)(x-x_0)}{r^{3}}\,\mathrm{d}V ,
\]
$F_y,F_z$ 只需把分子换成 $y-y_0$、$z-z_0$，其中 $r=\sqrt{(x-x_0)^2+(y-y_0)^2+(z-z_0)^2}$，$k$ 为引力常数。

\corepoint{分母是 $r^{\mathbf{3}}$ 而非 $r^2$：微元对单位质量的引力大小为 $\frac{k\mu\,\mathrm{d}V}{r^{2}}$，再乘方向余弦 $\frac{x-x_0}{r}$，合起来即 $\frac{x-x_0}{r^3}$。平面薄片对质点的引力把三重积分换成二重积分、$\mathrm{d}V$ 换成 $\mathrm{d}\sigma$，$r$ 仍按空间距离算。}

\section{含参变量的积分}

设 $f(x,y)$ 在矩形 $R=[a,b]\times[c,d]$ 上连续，则 $\Phi(y)=\int_a^bf(x,y)\,\mathrm{d}x$ 确定了 $y$ 在 $[c,d]$ 上的一个函数，称为\textbf{含参变量的积分}（$y$ 为参变量）。

\textbf{定理 1（连续性）}\quad 若 $f(x,y)$ 在 $R$ 上连续，则 $\Phi(y)$ 在 $[c,d]$ 上连续，即 $\lim\limits_{y\to y_0}\int_a^bf(x,y)\,\mathrm{d}x=\int_a^bf(x,y_0)\,\mathrm{d}x$（极限与积分可交换）。

\textbf{定理 2（积分次序交换）}\quad 若 $f$ 在 $R$ 上连续，则
\[
\int_c^d\mathrm{d}y\int_a^bf(x,y)\,\mathrm{d}x=\int_a^b\mathrm{d}x\int_c^df(x,y)\,\mathrm{d}y
\]
（即矩形区域上两种次序的二次积分相等）。

\textbf{定理 3（可导性·莱布尼茨公式）}\quad 若 $f$ 及其偏导数 $f_y$ 都在 $R$ 上连续，则 $\Phi(y)$ 在 $[c,d]$ 上可微，且
\keyformula{\Phi'(y)=\frac{\mathrm{d}}{\mathrm{d}y}\int_a^bf(x,y)\,\mathrm{d}x=\int_a^b\frac{\partial f(x,y)}{\partial y}\,\mathrm{d}x}
即\textbf{积分与求导可交换次序}。

\textbf{定理 4（积分限也含参变量）}\quad 若 $f$ 在 $R$ 上连续，$\alpha(y),\beta(y)$ 可微且 $a\le\alpha,\beta\le b$，则 $\Phi(y)=\int_{\alpha(y)}^{\beta(y)}f(x,y)\,\mathrm{d}x$ 可微，且
\[
\Phi'(y)=\int_{\alpha(y)}^{\beta(y)}f_y(x,y)\,\mathrm{d}x+f(\beta(y),y)\beta'(y)-f(\alpha(y),y)\alpha'(y).
\]

\textbf{常用方法}：(1) \textbf{对参变量微分法}——把待求定积分看作某含参积分的特殊值，用莱布尼茨公式求 $\Phi'(y)$，再积分回 $\Phi(y)$ 并代特殊值；(2) \textbf{先积分后求导}——先把 $\Phi(y)$ 算出再求导；(3) \textbf{交换积分次序}——把一元定积分凑成二重积分再换序。

\corepoint{用莱布尼茨公式前必须核对条件：$f$ 与 $f_y$ 都要在\textbf{闭矩形}上连续。积分限含参变量时别漏后两项。用"对参数微分法"求定积分时，引入参数要使 $\Phi(y_0)$ 好算，且全过程参数不得越出连续性所在的矩形。}
